\chapter[Constrained Optimization \& Convex Sets]{Constrained Optimization, Lagrange Multipliers, Convex Sets \& Convex Functions}
\label{chap:week09}

% ----------------------------------------------------------------------
% PART 1: EXECUTIVE ROADMAP & ARCHITECTURAL FOUNDATIONS
% ----------------------------------------------------------------------
\section{Executive Roadmap \& Learning Objectives}

\subsection{Opening Hook: The Boundary Fence and the Rolling Marble}
Imagine a vast, smooth hillside forming an elliptical bowl. At the very bottom of the bowl lies a deep depression: the unconstrained global minimum where a marble placed anywhere on the hillside would naturally roll and come to rest. 

Now imagine that a surveyor erects a rigid, straight wooden fence line cutting across the hillside high above the valley floor. If you release the marble high on the upper slope, where will it come to rest against the fence?

A common intuitive instinct is to look at where the bottom of the bowl lies, draw a straight perpendicular line from the valley floor to the fence, and assume the marble stops at that closest geometric point (the orthogonal Euclidean projection). But physical dynamics and mathematical optimization defy this naive projection! 

The marble does \emph{not} halt at the closest point to the valley floor. Gravity pulls the marble down the local slope; as it strikes the fence, the component of gravity parallel to the fence rail drives the marble rolling sideways along the wood. The marble continues to roll until the lateral downhill pull vanishes completely---that is, until the downward gravitational force is oriented purely perpendicular to the fence rail, exactly opposed by the normal contact force of the wood. 

The marble comes to rest at the precise point of \textbf{tangency}, where an elevation contour of the hill runs parallel to the fence line. 

In continuous mathematical optimization, this physical reality resolves the \textbf{Boundary Curse Paradox}. Finding a constrained minimum is fundamentally different from projecting an unconstrained minimum onto a constraint boundary. The contact force exerted by the fence is the physical embodiment of the \textbf{Lagrange Multiplier} $\nu^*$. It provides the exact normal counter-force required to balance the gradient of the objective functional.

\subsection{Architectural Bridge from Previous Chapters}
In Chapter~\ref{chap:week08}, we explored unconstrained optimization algorithms where parameter vectors $\mathbf{w}$ were permitted to roam freely across the entire Euclidean space $\mathbb{R}^d$. We analyzed gradient descent, Polyak momentum, and Newton-Raphson updates governed purely by local Taylor expansions and Hessian curvature.

However, in modern statistical machine learning, parameters are almost never unconstrained:
\begin{itemize}
    \item In multinomial classification, attention mechanisms, and Gaussian Mixture Models (Chapter~\ref{chap:week12}), parameter vectors represent probabilities and must reside on the \textbf{probability simplex}: $\sum_{i=1}^d w_i = 1$ with $w_i \ge 0$.
    \item In sparse feature selection via Lasso regression, parameters are bounded within an $L_1$ norm ball: $\|\mathbf{w}\|_1 \le s$.
    \item In hard-margin and soft-margin Support Vector Machines (Chapter~\ref{chap:week10}), classification margins enforce strict linear inequality boundaries: $y_i(\mathbf{w}^T \mathbf{x}_i + b) \ge 1$.
\end{itemize}
Chapter 9 bridges the gap between unconstrained vector calculus and constrained mathematical programming. We establish the geometric theory of Lagrange multipliers for manifolds, characterize convex sets and convex functions, prove the fundamental hyperplane separation theorems, and develop the mechanics of the Euclidean projection operator. This provides the mathematical launchpad for Karush-Kuhn-Tucker (KKT) duality in Chapter~\ref{chap:week10}.

\subsection{Prerequisites \& Scaffolding}
Mastery of this chapter builds upon vector calculus, linear geometry, and matrix analysis as organized in Table~\ref{tab:w09_prereqs}.

\begin{table}[htbp]
\centering
\small
\begin{tabularx}{\textwidth}{lXX}
\toprule
\textbf{Foundational Area} & \textbf{Core Concepts Required} & \textbf{Role in Machine Learning Foundations} \\
\midrule
Vector Calculus & Gradient vectors $\nabla f$, level surfaces $f(\mathbf{x}) = c$, tangent hyperplanes & Formulating orthogonality between objective gradients and constraint manifolds \\
\addlinespace
Linear Algebra & Hyperplanes $\mathbf{a}^T \mathbf{x} = b$, halfspaces $\mathbf{a}^T \mathbf{x} \le b$, orthogonal projections & Defining polyhedral feasible regions, normal cones, and projection operators \\
\addlinespace
Convex Analysis & Jensen's inequality, epigraphs, positive semidefinite cones $S^d_+$ & Establishing global optimality guarantees and non-expansiveness bounds \\
\bottomrule
\end{tabularx}
\caption{Prerequisite Scaffolding for Chapter 9.}
\label{tab:w09_prereqs}
\end{table}

\subsection{Core Competencies}
By the conclusion of this chapter, students will be able to:
\begin{itemize}
    \item Formulate equality-constrained optimization problems and apply the Method of Lagrange Multipliers.
    \item Prove that at a constrained local extremum, the objective gradient is collinear with the normal space of the constraint manifold.
    \item Define convex sets and verify domain convexity using intersection, affine map, and Cartesian product closure.
    \item State and apply the Separating Hyperplane Theorem and Supporting Hyperplane Theorem.
    \item Characterize convex functions via Jensen's inequality, first-order supporting hyperplanes, and second-order Hessian conditions.
    \item Prove the non-expansiveness of the Euclidean projection operator onto closed convex sets.
    \item Implement Projected Gradient Descent (PGD) with sparse $L_1$ ball projection.
\end{itemize}

\subsection{Notation \& Mathematical Conventions}
Throughout this chapter, we adhere to the standard mathematical conventions:
\begin{itemize}
    \item Feasible constraint set: $C \subseteq \mathbb{R}^d$.
    \item Euclidean projection operator: $\Pi_C(\mathbf{x}) \triangleq \argmin_{\mathbf{z} \in C} \|\mathbf{x} - \mathbf{z}\|_2$.
    \item Epigraph of a functional: $\operatorname{epi}(f) \triangleq \{(\mathbf{x}, t) \in \mathbb{R}^d \times \mathbb{R} \mid f(\mathbf{x}) \le t\}$.
    \item Sublevel set: $\mathcal{S}_\alpha(f) \triangleq \{\mathbf{x} \in \mathbb{R}^d \mid f(\mathbf{x}) \le \alpha\}$.
    \item Lagrangian functional: $\mathcal{L}(\mathbf{w}, \boldsymbol{\nu}) \triangleq f(\mathbf{w}) + \sum_{j=1}^p \nu_j h_j(\mathbf{w})$.
    \item Positive semidefinite cone: $S^d_+ \triangleq \{A \in \mathbb{R}^{d \times d} \mid A = A^T, \, \mathbf{v}^T A \mathbf{v} \ge 0 \; \forall \mathbf{v} \in \mathbb{R}^d\}$.
\end{itemize}

% ----------------------------------------------------------------------
% PART 2: MATHEMATICAL & COMPUTATIONAL FOUNDATIONS
% ----------------------------------------------------------------------
\section{Mathematical \& Computational Foundations}

\subsection{Equality Constraints and the Method of Lagrange Multipliers}
Consider the equality-constrained continuous optimization problem:
\begin{equation}
\min_{\mathbf{w} \in \mathbb{R}^d} f(\mathbf{w}) \quad \text{subject to } h_j(\mathbf{w}) = 0 \quad (j = 1, \dots, p)
\label{eq:eq_constrained_prob}
\end{equation}
where $f, h_1, \dots, h_p: \mathbb{R}^d \to \mathbb{R}$ are continuously differentiable functions with $p < d$. The feasible domain forms a $(d - p)$-dimensional geometric manifold:
\begin{equation}
\mathcal{M} = \{\mathbf{w} \in \mathbb{R}^d \mid h_j(\mathbf{w}) = 0, \; j = 1, \dots, p\}
\end{equation}

\begin{theorembox}{Geometric Condition of Optimality (Lagrange Multipliers)}
Let $\mathbf{w}^* \in \mathcal{M}$ be a local minimizer of $f$ on the feasible manifold $\mathcal{M}$. Assume the constraint gradients $\{\nabla h_1(\mathbf{w}^*), \dots, \nabla h_p(\mathbf{w}^*)\}$ are linearly independent (the \textbf{Linear Independence Constraint Qualification}, LICQ).

Then the gradient of the objective functional must reside strictly within the normal space spanned by the constraint gradients:
\begin{equation}
\nabla f(\mathbf{w}^*) + \sum_{j=1}^p \nu_j^* \nabla h_j(\mathbf{w}^*) = \mathbf{0}
\label{eq:lagrange_stationarity}
\end{equation}
for unique real scalars $\nu_1^*, \dots, \nu_p^* \in \mathbb{R}$ designated as \textbf{Lagrange Multipliers}.
\end{theorembox}

\begin{proof}
Let $\mathbf{w}^*$ be a local minimizer on $\mathcal{M}$. Consider any smooth curve $\boldsymbol{\gamma}: (-\epsilon, \epsilon) \to \mathcal{M}$ lying entirely within the feasible manifold such that $\boldsymbol{\gamma}(0) = \mathbf{w}^*$ and $\boldsymbol{\gamma}'(0) = \mathbf{v}$.
Since $\boldsymbol{\gamma}(t) \in \mathcal{M}$ for all $t$, every constraint satisfies $h_j(\boldsymbol{\gamma}(t)) = 0$ identically. Differentiating with respect to $t$ at $t = 0$ via the multivariable chain rule:
\begin{equation}
\frac{d}{dt} h_j(\boldsymbol{\gamma}(t))\Big|_{t=0} = \nabla h_j(\mathbf{w}^*)^T \boldsymbol{\gamma}'(0) = \nabla h_j(\mathbf{w}^*)^T \mathbf{v} = 0 \quad \forall j \in \{1, \dots, p\}
\end{equation}
Thus, the tangent space to the feasible manifold at $\mathbf{w}^*$ is precisely the orthogonal complement of the span of the constraint gradients:
\begin{equation}
T_{\mathcal{M}}(\mathbf{w}^*) = \left\{ \mathbf{v} \in \mathbb{R}^d \;\middle|\; \nabla h_j(\mathbf{w}^*)^T \mathbf{v} = 0, \, \forall j \right\} = \left( \Span\{\nabla h_1(\mathbf{w}^*), \dots, \nabla h_p(\mathbf{w}^*)\} \right)^\perp
\end{equation}
Now evaluate the rate of change of the objective functional along the feasible curve $\boldsymbol{\gamma}(t)$. Because $\mathbf{w}^*$ is a local minimum, the scalar function $\phi(t) = f(\boldsymbol{\gamma}(t))$ attains a local minimum at $t = 0$. By elementary calculus, $\phi'(0) = 0$:
\begin{equation}
\phi'(0) = \nabla f(\mathbf{w}^*)^T \boldsymbol{\gamma}'(0) = \nabla f(\mathbf{w}^*)^T \mathbf{v} = 0
\end{equation}
Because this orthogonality holds for \emph{every} feasible tangent vector $\mathbf{v} \in T_{\mathcal{M}}(\mathbf{w}^*)$:
\begin{equation}
\nabla f(\mathbf{w}^*) \in \left( T_{\mathcal{M}}(\mathbf{w}^*) \right)^\perp
\end{equation}
By the Fundamental Theorem of Linear Algebra (Chapter~\ref{chap:week02}), $\left( (V)^\perp \right)^\perp = V$. Therefore:
\begin{equation}
\nabla f(\mathbf{w}^*) \in \Span\{\nabla h_1(\mathbf{w}^*), \dots, \nabla h_p(\mathbf{w}^*)\}
\end{equation}
Expressing $\nabla f(\mathbf{w}^*)$ as a linear combination of $\{\nabla h_j(\mathbf{w}^*)\}$ with coefficients $-\nu_j^*$ yields~\eqref{eq:lagrange_stationarity}. Linear independence (LICQ) guarantees that these multiplier coefficients $\nu_j^*$ are uniquely determined.
\end{proof}

\begin{theorembox}{Lagrange Multiplier Sensitivity \& Shadow Price Theorem}
Let $f, h_1, \dots, h_p: \mathbb{R}^d \to \mathbb{R}$ be twice continuously differentiable. Consider the perturbed family of equality-constrained optimization problems:
\begin{equation}
f^*(\mathbf{c}) \triangleq \min_{\mathbf{w} \in \mathbb{R}^d} f(\mathbf{w}) \quad \text{subject to } h_j(\mathbf{w}) = c_j \quad (j = 1, \dots, p)
\label{eq:perturbed_opt_prob}
\end{equation}
where $\mathbf{c} = [c_1, \dots, c_p]^T \in \mathbb{R}^p$ represents a vector of constraint shifts.

Assume that at $\mathbf{c} = \mathbf{0}$, $\mathbf{w}^* = \mathbf{w}^*(\mathbf{0})$ is a local minimizer satisfying LICQ and the second-order sufficient condition ($\mathbf{v}^T \nabla_{\mathbf{w}}^2 \mathcal{L}(\mathbf{w}^*, \boldsymbol{\nu}^*) \mathbf{v} > 0$ for all non-zero $\mathbf{v} \in T_{\mathcal{M}}(\mathbf{w}^*)$).

Then the optimal value function $f^*(\mathbf{c})$ is continuously differentiable in an open neighborhood of $\mathbf{c} = \mathbf{0}$, and its gradient with respect to constraint perturbations is given exactly by the negative Lagrange multipliers:
\begin{equation}
\nabla_{\mathbf{c}} f^*(\mathbf{0}) = -\boldsymbol{\nu}^* \iff \frac{\partial f^*(\mathbf{0})}{\partial c_j} = -\nu_j^* \quad (j = 1, \dots, p)
\label{eq:lagrange_sensitivity_grad}
\end{equation}
\end{theorembox}

\begin{proof}
By the Implicit Function Theorem applied to the stationary Karush-Kuhn-Tucker system:
\begin{equation}
F(\mathbf{w}, \boldsymbol{\nu}, \mathbf{c}) = \begin{bmatrix} \nabla f(\mathbf{w}) + \sum_{j=1}^p \nu_j \nabla h_j(\mathbf{w}) \\ h_1(\mathbf{w}) - c_1 \\ \vdots \\ h_p(\mathbf{w}) - c_p \end{bmatrix} = \mathbf{0}
\end{equation}
the Jacobian matrix with respect to $(\mathbf{w}, \boldsymbol{\nu})$ is:
\begin{equation}
J_{(\mathbf{w}, \boldsymbol{\nu})} = \begin{bmatrix} \nabla_{\mathbf{w}}^2 \mathcal{L}(\mathbf{w}, \boldsymbol{\nu}) & \nabla \mathbf{h}(\mathbf{w}) \\ \nabla \mathbf{h}(\mathbf{w})^T & O \end{bmatrix}
\end{equation}
Under LICQ and the second-order sufficient condition, this bordered Hessian matrix is non-singular. Consequently, there exist unique, continuously differentiable local solution trajectories $\mathbf{w}^*(\mathbf{c})$ and $\boldsymbol{\nu}^*(\mathbf{c})$ satisfying $\mathbf{w}^*(\mathbf{0}) = \mathbf{w}^*$, $\boldsymbol{\nu}^*(\mathbf{0}) = \boldsymbol{\nu}^*$, and $h_j(\mathbf{w}^*(\mathbf{c})) = c_j$ identically.

Now evaluate the derivative of the optimal value function $f^*(\mathbf{c}) = f(\mathbf{w}^*(\mathbf{c}))$ with respect to coordinate $c_k$ via the multivariable chain rule:
\begin{equation}
\frac{\partial f^*(\mathbf{c})}{\partial c_k} = \sum_{i=1}^d \frac{\partial f(\mathbf{w}^*(\mathbf{c}))}{\partial w_i} \frac{\partial w_i^*(\mathbf{c})}{\partial c_k} = \nabla f(\mathbf{w}^*(\mathbf{c}))^T \frac{\partial \mathbf{w}^*(\mathbf{c})}{\partial c_k}
\end{equation}
Substitute the stationarity condition $\nabla f(\mathbf{w}^*(\mathbf{c})) = -\sum_{j=1}^p \nu_j^*(\mathbf{c}) \nabla h_j(\mathbf{w}^*(\mathbf{c}))$:
\begin{equation}
\frac{\partial f^*(\mathbf{c})}{\partial c_k} = -\sum_{j=1}^p \nu_j^*(\mathbf{c}) \left( \nabla h_j(\mathbf{w}^*(\mathbf{c}))^T \frac{\partial \mathbf{w}^*(\mathbf{c})}{\partial c_k} \right)
\label{eq:chain_rule_sub}
\end{equation}
Next, differentiate the identity $h_j(\mathbf{w}^*(\mathbf{c})) = c_j$ with respect to $c_k$ on both sides:
\begin{equation}
\frac{\partial}{\partial c_k} h_j(\mathbf{w}^*(\mathbf{c})) = \nabla h_j(\mathbf{w}^*(\mathbf{c}))^T \frac{\partial \mathbf{w}^*(\mathbf{c})}{\partial c_k} = \frac{\partial c_j}{\partial c_k} = \delta_{jk}
\end{equation}
where $\delta_{jk}$ is the Kronecker delta ($\delta_{jk} = 1$ if $j = k$, and $0$ otherwise).

Substituting this identity into~\eqref{eq:chain_rule_sub}:
\begin{equation}
\frac{\partial f^*(\mathbf{c})}{\partial c_k} = -\sum_{j=1}^p \nu_j^*(\mathbf{c}) \delta_{jk} = -\nu_k^*(\mathbf{c})
\end{equation}
Evaluating at $\mathbf{c} = \mathbf{0}$ yields the fundamental shadow price identity $\frac{\partial f^*(\mathbf{0})}{\partial c_k} = -\nu_k^*$.
\end{proof}

\begin{ponderbox}{What Do Lagrange Multipliers Actually Measure?}
What is the physical and economic significance of the optimal Lagrange multiplier $\nu_j^*$?
\tcblower
\textbf{The Socratic Resolution:} The Lagrange multiplier measures the \textbf{marginal sensitivity (shadow price)} of the optimal objective value with respect to perturbations of the constraint boundary. 

Suppose the constraint is perturbed by a small scalar shift: $h_j(\mathbf{w}) = \epsilon$. Let $\mathbf{w}^*(\epsilon)$ denote the perturbed optimal solution, and let $f^*(\epsilon) = f(\mathbf{w}^*(\epsilon))$ denote the perturbed optimal loss. Applying the chain rule:
\begin{equation}
\frac{d f^*(\epsilon)}{d\epsilon}\Big|_{\epsilon = 0} = \nabla f(\mathbf{w}^*)^T \frac{d\mathbf{w}^*}{d\epsilon}
\end{equation}
Differentiating the perturbed constraint $h_j(\mathbf{w}^*(\epsilon)) = \epsilon$ yields $\nabla h_j(\mathbf{w}^*)^T \frac{d\mathbf{w}^*}{d\epsilon} = 1$. Substituting the Lagrange stationarity condition $\nabla f(\mathbf{w}^*) = -\sum_k \nu_k^* \nabla h_k(\mathbf{w}^*)$:
\begin{equation}
\frac{d f^*(\epsilon)}{d\epsilon}\Big|_{\epsilon = 0} = -\sum_{k=1}^p \nu_k^* \nabla h_k(\mathbf{w}^*)^T \frac{d\mathbf{w}^*}{d\epsilon} = -\nu_j^*(1) = -\nu_j^*
\end{equation}
Thus, $-\nu_j^*$ is the exact marginal derivative of optimal loss with respect to constraint tightening. If $|\nu_j^*|$ is massive, shifting the constraint boundary slightly will severely degrade performance. If $\nu_j^* = 0$, the constraint has zero marginal impact on the optimal objective value.
\end{ponderbox}

\subsection{Convex Sets and Hyperplane Separation}

\begin{definitionbox}{Convex Set}
A subset $C \subseteq \mathbb{R}^d$ is \textbf{convex} if the line segment connecting any two arbitrary points in $C$ lies entirely within $C$:
\begin{equation}
\forall \mathbf{x}, \mathbf{y} \in C, \; \forall \theta \in [0, 1] \implies \theta \mathbf{x} + (1 - \theta)\mathbf{y} \in C
\label{eq:convex_set_def}
\end{equation}
A point of the form $\sum_{i=1}^k \theta_i \mathbf{x}_i$ with $\theta_i \ge 0$ and $\sum_{i=1}^k \theta_i = 1$ is a \textbf{convex combination}. A set is convex if and only if it contains all convex combinations of its elements.
\end{definitionbox}

\begin{theorembox}{Operations Preserving Set Convexity}
Convexity of sets is invariant under the following fundamental mathematical operations:
\begin{enumerate}
    \item \textbf{Arbitrary Intersections:} If $C_\alpha$ is a convex set for every index $\alpha \in \mathcal{A}$ (including uncountably infinite index collections), then their intersection is convex:
    \begin{equation}
    C = \bigcap_{\alpha \in \mathcal{A}} C_\alpha \quad \text{is convex.}
    \end{equation}
    Consequently, polyhedra defined as intersections of finitely many halfspaces $\{\mathbf{x} \mid A\mathbf{x} \le \mathbf{b}\}$ are always convex sets.
    \item \textbf{Affine Transformations:} If $C \subseteq \mathbb{R}^n$ is convex and $f: \mathbb{R}^n \to \mathbb{R}^m$ is an affine map $f(\mathbf{x}) = A\mathbf{x} + \mathbf{b}$, then the image $f(C) = \{A\mathbf{x} + \mathbf{b} \mid \mathbf{x} \in C\}$ is convex. Furthermore, the inverse image (pre-image) $f^{-1}(C) = \{\mathbf{x} \mid A\mathbf{x} + \mathbf{b} \in C\}$ is convex.
    \item \textbf{Cartesian Products:} If $C_1 \subseteq \mathbb{R}^n$ and $C_2 \subseteq \mathbb{R}^m$ are convex sets, their Cartesian product:
    \begin{equation}
    C_1 \times C_2 = \{(\mathbf{x}_1, \mathbf{x}_2) \mid \mathbf{x}_1 \in C_1, \; \mathbf{x}_2 \in C_2\} \subseteq \mathbb{R}^{n+m} \quad \text{is convex.}
    \end{equation}
\end{enumerate}
\end{theorembox}

\begin{theorembox}{Separating Hyperplane Theorem}
Let $C_1, C_2 \subset \mathbb{R}^d$ be two non-empty, disjoint convex sets ($C_1 \cap C_2 = \emptyset$). 
Then there exists a non-zero normal vector $\mathbf{a} \in \mathbb{R}^d \setminus \{\mathbf{0}\}$ and a scalar threshold $b \in \mathbb{R}$ such that:
\begin{equation}
\mathbf{a}^T \mathbf{x} \le b \le \mathbf{a}^T \mathbf{y} \quad \forall \mathbf{x} \in C_1, \; \forall \mathbf{y} \in C_2
\label{eq:separating_hyperplane}
\end{equation}
The affine hyperplane $\mathcal{H} = \{\mathbf{z} \in \mathbb{R}^d \mid \mathbf{a}^T \mathbf{z} = b\}$ separates $C_1$ and $C_2$. If $C_1$ is closed and $C_2$ is compact (or a singleton outside $C_1$), the separation is \textbf{strict}: $\mathbf{a}^T \mathbf{x} < b < \mathbf{a}^T \mathbf{y}$.
\end{theorembox}

\begin{ponderbox}{Why Intersection Preserves Convexity While Union Destroys It}
Why does an intersection of arbitrarily many convex sets always remain convex, while the union of even two convex sets almost never remains convex?
\tcblower
\textbf{The Socratic Resolution:} Convexity is an \textbf{internal closure property} across lines. 
If $\mathbf{x}, \mathbf{y} \in \bigcap_{\alpha} C_\alpha$, then $\mathbf{x}$ and $\mathbf{y}$ belong to every single set $C_\alpha$. Because each individual $C_\alpha$ is convex, the entire line segment $[\mathbf{x}, \mathbf{y}] = \{\theta \mathbf{x} + (1-\theta)\mathbf{y} \mid \theta \in [0, 1]\}$ lies inside $C_\alpha$ for every $\alpha$. Therefore, the entire segment lies in the intersection: $[\mathbf{x}, \mathbf{y}] \subseteq \bigcap_\alpha C_\alpha$.

Conversely, consider the union of two disjoint convex sets $C_1 \cup C_2$. Pick $\mathbf{x} \in C_1$ and $\mathbf{y} \in C_2$. The line segment connecting $\mathbf{x}$ to $\mathbf{y}$ must physically traverse the empty void separating the two sets. Points along the segment belong to neither $C_1$ nor $C_2$, instantly violating the convexity condition.
\end{ponderbox}

\subsection{Convex Functions}

\begin{definitionbox}{Convex Function}
A function $f: C \to \mathbb{R}$ defined on a convex domain $C \subseteq \mathbb{R}^d$ is \textbf{convex} if for all $\mathbf{x}, \mathbf{y} \in C$ and all $\theta \in [0, 1]$:
\begin{equation}
f(\theta \mathbf{x} + (1 - \theta)\mathbf{y}) \le \theta f(\mathbf{x}) + (1 - \theta) f(\mathbf{y})
\label{eq:convex_func_def}
\end{equation}
A function is \textbf{strictly convex} if the inequality is strict for all $\mathbf{x} \ne \mathbf{y}$ and $\theta \in (0, 1)$. Geometrically, convexity states that the secant chord connecting $( \mathbf{x}, f(\mathbf{x}) )$ and $( \mathbf{y}, f(\mathbf{y}) )$ lies everywhere on or above the graph of the function.
\end{definitionbox}

\begin{definitionbox}{Epigraph Characterization}
The \textbf{epigraph} of a function $f: C \to \mathbb{R}$ is the subset of $\mathbb{R}^{d+1}$ lying on or above its graph:
\begin{equation}
\operatorname{epi}(f) \triangleq \left\{ (\mathbf{x}, t) \in \mathbb{R}^d \times \mathbb{R} \;\middle|\; \mathbf{x} \in C, \; f(\mathbf{x}) \le t \right\}
\label{eq:epigraph_def}
\end{equation}
\textbf{The Fundamental Bridge:} A function $f$ is convex if and only if its epigraph $\operatorname{epi}(f)$ is a convex set in $\mathbb{R}^{d+1}$.
\end{definitionbox}

\begin{theorembox}{First-Order Characterization of Convexity}
Let $f: C \to \mathbb{R}$ be continuously differentiable on an open convex set $C$. Then $f$ is convex if and only if for all $\mathbf{x}, \mathbf{y} \in C$:
\begin{equation}
f(\mathbf{y}) \ge f(\mathbf{x}) + \nabla f(\mathbf{x})^T (\mathbf{y} - \mathbf{x})
\label{eq:first_order_convex}
\end{equation}
Geometrically, the first-order Taylor approximation forms a global supporting hyperplane that lies everywhere below or touches the graph of $f$.
\end{theorembox}

\begin{proof}
Assume $f$ is convex. For any $\theta \in (0, 1]$:
\begin{equation}
f(\mathbf{x} + \theta(\mathbf{y} - \mathbf{x})) = f((1 - \theta)\mathbf{x} + \theta \mathbf{y}) \le (1 - \theta)f(\mathbf{x}) + \theta f(\mathbf{y})
\end{equation}
Subtracting $f(\mathbf{x})$ and dividing by $\theta$:
\begin{equation}
\frac{f(\mathbf{x} + \theta(\mathbf{y} - \mathbf{x})) - f(\mathbf{x})}{\theta} \le f(\mathbf{y}) - f(\mathbf{x})
\end{equation}
Taking the limit as $\theta \to 0^+$, the left-hand side converges to the directional derivative $\nabla f(\mathbf{x})^T (\mathbf{y} - \mathbf{x})$, yielding $f(\mathbf{y}) \ge f(\mathbf{x}) + \nabla f(\mathbf{x})^T (\mathbf{y} - \mathbf{x})$.

Conversely, suppose~\eqref{eq:first_order_convex} holds for all points. Let $\mathbf{z} = \theta \mathbf{x} + (1 - \theta)\mathbf{y}$ for $\theta \in [0, 1]$. Applying~\eqref{eq:first_order_convex} at $\mathbf{z}$ with respect to $\mathbf{x}$ and $\mathbf{y}$:
\begin{align}
f(\mathbf{x}) &\ge f(\mathbf{z}) + \nabla f(\mathbf{z})^T (\mathbf{x} - \mathbf{z}) \\
f(\mathbf{y}) &\ge f(\mathbf{z}) + \nabla f(\mathbf{z})^T (\mathbf{y} - \mathbf{z})
\end{align}
Multiplying the first inequality by $\theta$, the second by $(1 - \theta)$, and adding:
\begin{align}
\theta f(\mathbf{x}) + (1 - \theta)f(\mathbf{y}) &\ge f(\mathbf{z}) + \nabla f(\mathbf{z})^T \left( \theta(\mathbf{x} - \mathbf{z}) + (1 - \theta)(\mathbf{y} - \mathbf{z}) \right) \nonumber \\
&= f(\mathbf{z}) + \nabla f(\mathbf{z})^T (\mathbf{0}) = f(\theta \mathbf{x} + (1 - \theta)\mathbf{y})
\end{align}
which proves that $f$ is convex.
\end{proof}

\begin{theorembox}{Second-Order Characterization of Convexity}
Let $f: C \to \mathbb{R}$ be twice continuously differentiable on an open convex set $C \subseteq \mathbb{R}^d$. Then:
\begin{itemize}
    \item $f$ is convex if and only if its Hessian matrix is positive semidefinite everywhere:
    \begin{equation}
    \nabla^2 f(\mathbf{x}) \succeq 0 \quad \forall \mathbf{x} \in C
    \end{equation}
    \item If the Hessian is strictly positive definite everywhere ($\nabla^2 f(\mathbf{x}) \succ 0$ for all $\mathbf{x} \in C$), then $f$ is \textbf{strictly convex}.
\end{itemize}
\end{theorembox}

\begin{theorembox}{The Strong Convexity \& Smoothness Quadratic Sandwich}
Let $f: \mathbb{R}^d \to \mathbb{R}$ be twice continuously differentiable.
\begin{enumerate}
    \item $f$ is \textbf{$m$-strongly convex} ($m > 0$) if and only if $\nabla^2 f(\mathbf{x}) \succeq m I$ for all $\mathbf{x} \in \mathbb{R}^d$. This is equivalent to the global quadratic under-estimator:
    \begin{equation}
    f(\mathbf{y}) \ge f(\mathbf{x}) + \nabla f(\mathbf{x})^T(\mathbf{y} - \mathbf{x}) + \frac{m}{2}\|\mathbf{y} - \mathbf{x}\|_2^2 \quad \forall \mathbf{x}, \mathbf{y} \in \mathbb{R}^d
    \label{eq:strong_convexity_lower}
    \end{equation}
    \item $f$ is \textbf{$L$-smooth} ($L > 0$) if and only if its gradient is $L$-Lipschitz continuous ($\|\nabla f(\mathbf{x}) - \nabla f(\mathbf{y})\|_2 \le L \|\mathbf{x} - \mathbf{y}\|_2$), which is equivalent to $\nabla^2 f(\mathbf{x}) \preceq L I$ for all $\mathbf{x} \in \mathbb{R}^d$. This implies the global quadratic over-estimator (Descent Lemma):
    \begin{equation}
    f(\mathbf{y}) \le f(\mathbf{x}) + \nabla f(\mathbf{x})^T(\mathbf{y} - \mathbf{x}) + \frac{L}{2}\|\mathbf{y} - \mathbf{x}\|_2^2 \quad \forall \mathbf{x}, \mathbf{y} \in \mathbb{R}^d
    \label{eq:smoothness_upper}
    \end{equation}
\end{enumerate}
Combining both bounds, any $m$-strongly convex and $L$-smooth functional is strictly sandwiched between two quadratic paraboloids of curvature $m$ and $L$:
\begin{equation}
f(\mathbf{x}) + \nabla f(\mathbf{x})^T(\mathbf{y} - \mathbf{x}) + \frac{m}{2}\|\mathbf{y} - \mathbf{x}\|_2^2 \le f(\mathbf{y}) \le f(\mathbf{x}) + \nabla f(\mathbf{x})^T(\mathbf{y} - \mathbf{x}) + \frac{L}{2}\|\mathbf{y} - \mathbf{x}\|_2^2
\label{eq:quadratic_sandwich}
\end{equation}
The ratio $\kappa \triangleq \frac{L}{m} \ge 1$ is the \textbf{condition number} of the objective functional, dictating the geometric eccentricity of all sublevel sets and governing the linear convergence rate $(1 - 1/\kappa)$ of both Gradient Descent and Projected Gradient Descent.
\end{theorembox}

\begin{proof}
By Taylor's theorem with integral remainder applied to the scalar function $g(t) = f(\mathbf{x} + t(\mathbf{y} - \mathbf{x}))$ for $t \in [0, 1]$:
\begin{equation}
f(\mathbf{y}) = f(\mathbf{x}) + \nabla f(\mathbf{x})^T(\mathbf{y} - \mathbf{x}) + \int_0^1 (1 - t) (\mathbf{y} - \mathbf{x})^T \nabla^2 f(\mathbf{x} + t(\mathbf{y} - \mathbf{x})) (\mathbf{y} - \mathbf{x}) \, dt
\end{equation}
If $\nabla^2 f(\mathbf{z}) \succeq m I$ for all $\mathbf{z}$, then $(\mathbf{y} - \mathbf{x})^T \nabla^2 f(\mathbf{z}) (\mathbf{y} - \mathbf{x}) \ge m \|\mathbf{y} - \mathbf{x}\|_2^2$. Integrating:
\begin{equation}
\int_0^1 (1 - t) m \|\mathbf{y} - \mathbf{x}\|_2^2 \, dt = m \|\mathbf{y} - \mathbf{x}\|_2^2 \left[ t - \frac{t^2}{2} \right]_0^1 = \frac{m}{2} \|\mathbf{y} - \mathbf{x}\|_2^2
\end{equation}
which proves~\eqref{eq:strong_convexity_lower}. Applying the identical integration to the upper bound $(\mathbf{y} - \mathbf{x})^T \nabla^2 f(\mathbf{z}) (\mathbf{y} - \mathbf{x}) \le L \|\mathbf{y} - \mathbf{x}\|_2^2$ yields~\eqref{eq:smoothness_upper}.
\end{proof}

\begin{ponderbox}{Convex Sets vs.\ Convex Functions: The Quasiconvexity Ambiguity}
Every convex function has convex sublevel sets $\mathcal{S}_\alpha = \{\mathbf{x} \mid f(\mathbf{x}) \le \alpha\}$. Does having convex sublevel sets guarantee that a function is convex?
\tcblower
\textbf{The Socratic Resolution:} No! Functions whose sublevel sets are all convex are called \textbf{quasiconvex}, and quasiconvexity is strictly weaker than convexity. 

Consider the one-dimensional function $f(x) = \sqrt{|x|}$. For any $\alpha \ge 0$, the sublevel set is $\mathcal{S}_\alpha = [-\alpha^2, \alpha^2]$, which is an interval and therefore a convex set. Yet $f(x)$ is blatantly non-convex: between $x = 0$ and $x = 4$, the secant chord lies \emph{below} the function curve ($f(2) = \sqrt{2} \approx 1.414$, while the chord midpoint is $\frac{1}{2}(0 + 2) = 1.0 < 1.414$).

Similarly, the Gaussian probability density function $f(x) = e^{-x^2}$ has convex upper level sets (it is quasiconcave), but fails standard concavity. Convexity requires the epigraph in $\mathbb{R}^{d+1}$ to be convex; having convex horizontal cross-sections (sublevel sets) does not ensure that the vertical profile satisfies the secant inequality.
\end{ponderbox}

% ----------------------------------------------------------------------
% PART 3: VISUAL & GEOMETRIC ANCHORING
% ----------------------------------------------------------------------
\section{Visual \& Geometric Anchoring}

Optimization geometry clarifies mathematical mechanics. In this section, we anchor the theoretical conditions in three native TikZ visualizations.

\subsection{The Geometry of Lagrange Multipliers and Tangency}
Figure~\ref{fig:constrained_geometry} contrasts the true constrained minimum with the naive orthogonal projection of the unconstrained minimum, demonstrating why constrained extrema are points of contour tangency.

\begin{figure}[htbp]
\centering
\begin{tikzpicture}[scale=1.0, font=\sffamily\footnotesize]
    % Axes
    \draw[->, thick, color=primarydark!50] (-0.8, 0) -- (4.2, 0) node[right] {$x$};
    \draw[->, thick, color=primarydark!50] (0, -0.8) -- (0, 4.2) node[above] {$y$};

    % Objective Contours: f(x, y) = 2x^2 + y^2 = c => x^2 / (c/2) + y^2 / c = 1
    % c = 2: rx = 1.0, ry = 1.414
    \draw[thick, color=primarydark!35] (0, 0) ellipse (1.0cm and 1.414cm);
    \node[right, color=primarydark!40] at (1.0, 0.4) {\scriptsize $f = 2$};

    % c = 6: rx = 1.732, ry = 2.449 (TANGENT CONTOUR!)
    \draw[very thick, color=primarydark!85] (0, 0) ellipse (1.732cm and 2.449cm);
    \node[right, color=primarydark!90] at (1.75, 1.2) {\textbf{Optimal Contour} $f = 6.0$};

    % c = 6.75: rx = 1.837, ry = 2.598 (SUBOPTIMAL PROJECTION CONTOUR)
    \draw[thick, dashed, color=traprad!65] (0, 0) ellipse (1.837cm and 2.598cm);
    \node[right, color=traprad!80] at (1.85, 2.1) {\scriptsize Naive Contour $f = 6.75$};

    % Constraint Line: x + y = 3
    \draw[line width=1.8pt, color=tealborder] (-0.3, 3.3) -- (3.3, -0.3);
    \node[above right, color=tealborder] at (2.4, 0.7) {\textbf{Feasible Constraint Line } $x + y = 3$};

    % Unconstrained minimum
    \filldraw[color=primarydark] (0, 0) circle (2.5pt) node[below left] {$\mathbf{w}_{\text{unc}}^* = (0, 0)$};

    % Naive Euclidean Projection: (1.5, 1.5)
    \coordinate (Pnaive) at (1.5, 1.5);
    \filldraw[color=traprad] (Pnaive) circle (3pt);
    \draw[dashed, color=traprad, thick] (0, 0) -- (Pnaive);
    \node[above right, color=traprad, align=left] at (1.55, 1.55) {
        \textbf{Naive Projection } $\Pi_C(\mathbf{0}) = (1.5, 1.5)$ \\
        \scriptsize $f = 6.75$ ($+12.5\%$ suboptimality)
    };

    % True Constrained Minimum: (1.0, 2.0)
    \coordinate (Wstar) at (1.0, 2.0);
    \filldraw[color=forest] (Wstar) circle (3.5pt);
    \node[above left, color=forest, align=right] at (0.95, 2.05) {
        \textbf{True Constrained Min } $\mathbf{w}_C^* = (1, 2)$ \\
        \scriptsize $f^* = 6.00$ (\textbf{Exact Tangency})
    };

    % Collinear gradient arrows at (1, 2)
    % grad f(1, 2) = [4, 4]^T => scale down by 0.25 => [1.0, 1.0]
    \draw[->, >=stealth, line width=2pt, color=coralborder] (Wstar) -- ($(Wstar) + (0.8, 0.8)$) 
        node[right] {$\nabla f = [4, 4]^T$};
    % grad h(1, 2) = [1, 1]^T => scale down by 0.6 => [0.6, 0.6]
    \draw[->, >=stealth, line width=1.6pt, color=tealborder] (Wstar) -- ($(Wstar) + (-0.6, -0.6)$) 
        node[below left] {$\nu^* \nabla h = -4 [1, 1]^T$};
\end{tikzpicture}
\caption{The Geometry of Lagrange Multipliers: Curvature dictates that constrained optima are points of contour tangency, not orthogonal projections. The unconstrained minimum is at $(0, 0)$. The naive Euclidean projection of $(0, 0)$ onto the constraint line $x + y = 3$ is $(1.5, 1.5)$, yielding a suboptimal loss of $f = 6.75$. The true constrained minimum occurs at $(1, 2)$, where the objective ellipse $f = 6.00$ is tangent to the line, aligning the gradient vectors $\nabla f$ and $\nabla h$ co-linearly.}
\label{fig:constrained_geometry}
\end{figure}

\subsection{Convex Sets and the Separating Hyperplane Theorem}
Figure~\ref{fig:convex_sets_separation} diagrams the internal chord test for set convexity and illustrates the geometric separation of disjoint convex bodies by an affine hyperplane.

\begin{figure}[htbp]
\centering
\begin{tikzpicture}[scale=0.9, font=\sffamily\footnotesize]
    % Panel 1: Convex Set
    \begin{scope}[shift={(-5.0, 0)}]
        \draw[thick, fill=tealborder!15, color=tealborder] 
            (-1.5, -1.0) to[out=90, in=180] (0, 1.5)
                         to[out=0, in=90] (1.8, 0.2)
                         to[out=-90, in=0] (0.8, -1.5)
                         to[out=180, in=-90] cycle;
        \coordinate (X1) at (-0.8, -0.3);
        \coordinate (Y1) at (1.0, 0.6);
        \draw[very thick, color=forest] (X1) -- (Y1);
        \filldraw[color=forest] (X1) circle (2.5pt) node[below left] {$\mathbf{x}$};
        \filldraw[color=forest] (Y1) circle (2.5pt) node[above right] {$\mathbf{y}$};
        \node[above, color=forest] at (0.1, 0.2) {\scriptsize $[\mathbf{x}, \mathbf{y}] \subset C$};
        \node[below, align=center] at (0, -2.0) {\textbf{(a) Convex Set} \\ All secant segments contained};
    \end{scope}

    % Panel 2: Non-Convex Set (Crescent)
    \begin{scope}[shift={(0, 0)}]
        \draw[thick, fill=traprad!15, color=traprad] 
            (-1.5, -1.2) to[out=80, in=180] (0, 1.5)
                         to[out=0, in=100] (1.5, -1.2)
                         to[out=120, in=0] (0, 0.2)
                         to[out=180, in=60] cycle;
        \coordinate (X2) at (-1.0, 0.2);
        \coordinate (Y2) at (1.0, 0.2);
        \draw[very thick, dashed, color=traprad] (X2) -- (Y2);
        \filldraw[color=traprad] (X2) circle (2.5pt) node[left] {$\mathbf{x}$};
        \filldraw[color=traprad] (Y2) circle (2.5pt) node[right] {$\mathbf{y}$};
        \node[above, color=traprad] at (0, 0.3) {\scriptsize Segment escapes!};
        \node[below, align=center] at (0, -2.0) {\textbf{(b) Non-Convex Set} \\ Segment exits the set};
    \end{scope}

    % Panel 3: Separating Hyperplane
    \begin{scope}[shift={(5.2, 0)}]
        % Set C1
        \draw[thick, fill=primarydark!15, color=primarydark] 
            (-1.8, -0.8) circle (0.9cm);
        \node at (-1.8, -0.8) {$C_1$};

        % Set C2
        \draw[thick, fill=forest!15, color=forest] 
            (1.5, 0.5) ellipse (1.1cm and 0.7cm);
        \node at (1.5, 0.5) {$C_2$};

        % Separating Hyperplane line
        \draw[very thick, color=coralborder] (-0.3, -2.0) -- (0.5, 2.0) 
            node[above] {\textbf{Hyperplane } $\mathbf{a}^T\mathbf{x} = b$};
        
        % Normal vector a
        \draw[->, >=stealth, line width=1.5pt, color=coralborder] (0.1, 0) -- (1.1, -0.2) 
            node[below right] {$\mathbf{a}$};

        \node[below, align=center] at (0, -2.0) {\textbf{(c) Separating Hyperplane} \\ Disjoint convex sets};
    \end{scope}
\end{tikzpicture}
\caption{Convex Geometry and the Separating Hyperplane Theorem: (a) In a convex set, the straight segment between any two points $\mathbf{x}, \mathbf{y}$ remains entirely enclosed within the set. (b) In a non-convex set (e.g., an indented crescent), chords cross through the external void. (c) The Separating Hyperplane Theorem guarantees that any two non-empty disjoint convex sets $C_1, C_2$ can be completely separated by an affine hyperplane $\mathbf{a}^T\mathbf{x} = b$.}
\label{fig:convex_sets_separation}
\end{figure}

\subsection{First-Order Supporting Hyperplanes and Epigraph Geometry}
Figure~\ref{fig:epigraph_supporting} illustrates the equivalence between function convexity, epigraph convexity, and the global supporting hyperplane property.

\begin{figure}[htbp]
\centering
\begin{tikzpicture}[scale=1.1, font=\sffamily\footnotesize]
    % Axes
    \draw[->, thick, color=primarydark!60] (-3.2, 0) -- (3.2, 0) node[right] {$x$};
    \draw[->, thick, color=primarydark!60] (0, -0.5) -- (0, 4.5) node[above] {$t$};

    % Shaded Epigraph region: t >= 0.4*x^2 + 0.8
    \fill[tealborder!15, domain=-2.8:2.8] 
        (-2.8, 4.2) -- plot (\x, {0.4*\x*\x + 0.8}) -- (2.8, 4.2) -- cycle;
    \node[color=tealborder!90!black] at (1.4, 3.2) {\large \textbf{Epigraph } $\operatorname{epi}(f)$};

    % Convex function curve f(x) = 0.4*x^2 + 0.8
    \draw[line width=1.8pt, color=primarydark, domain=-2.8:2.8] 
        plot (\x, {0.4*\x*\x + 0.8}) node[right] {$t = f(x)$};

    % Base point x0 = -1.0 => f(-1.0) = 0.4(1) + 0.8 = 1.2
    % f'(x) = 0.8*x => f'(-1.0) = -0.8
    % Tangent line: t = 1.2 - 0.8*(x - (-1.0)) = 1.2 - 0.8(x + 1) = 0.4 - 0.8*x
    \coordinate (X0) at (-1.0, 1.2);
    \filldraw[color=forest] (X0) circle (2.5pt) node[above left=2pt] {$(x_0, f(x_0))$};

    % Supporting Hyperplane (Tangent line)
    \draw[very thick, color=coralborder, domain=-3.0:2.5] 
        plot (\x, {1.2 - 0.8*(\x + 1.0)}) 
        node[below right] {\textbf{Supporting Hyperplane} $t = f(x_0) + f'(x_0)(x - x_0)$};

    % Normal vector to supporting hyperplane: [f'(x0), -1]^T = [-0.8, -1]^T pointing downward
    % or [-f'(x0), 1]^T pointing into epigraph
    \draw[->, >=stealth, line width=1.5pt, color=forest] (X0) -- ($(X0) + (0.8, 1.0)$) 
        node[above] {\scriptsize Inward normal $(-\nabla f(x_0), 1)$};

    % Annotation demonstrating global under-estimation
    \draw[<->, >=stealth, thick, color=traprad] (1.5, {0.4*(1.5)*(1.5) + 0.8}) -- (1.5, {1.2 - 0.8*(1.5 + 1.0)});
    \node[right, color=traprad, align=left] at (1.55, 0.4) {
        \scriptsize Gap $\ge 0$ \\
        \scriptsize $f(y) \ge f(x_0) + \nabla f(x_0)^T(y - x_0)$
    };
\end{tikzpicture}
\caption{First-Order Characterization and Epigraph Geometry: A functional $f$ is convex if and only if its epigraph $\operatorname{epi}(f)$ is a convex set in $\mathbb{R}^{d+1}$. Consequently, at every point $x_0$, the affine first-order Taylor expansion defines a global supporting hyperplane that under-estimates the functional everywhere across its domain.}
\label{fig:epigraph_supporting}
\end{figure}

\subsection{The Geometry of Water-Filling and Resource Allocation}
Figure~\ref{fig:waterfilling_geometry} visualizes the fundamental \textbf{Water-Filling Principle} originating in Shannon information theory and convex resource allocation. When allocating a finite total resource (such as transmit power $P_{\text{total}}$ or portfolio capital) across independent channels subjected to varying noise levels $\sigma_i^2$, the optimal Lagrange multiplier establishes a global equilibrium threshold $\mu^* \triangleq 1/\nu^*$. 

\begin{figure}[htbp]
\centering
\begin{tikzpicture}[scale=1.0, font=\sffamily\footnotesize]
    % Channel 1 basin: x in [0.5, 2.5], floor at y = 0.4
    \draw[thick, fill=gray!20] (0.5, 0.4) rectangle (2.5, 0);
    \draw[thick] (0.5, 3.5) -- (0.5, 0.4) -- (2.5, 0.4) -- (2.5, 3.5);
    \node at (1.5, 0.2) {\textbf{Noise } $\sigma_1^2 = 0.2$};
    
    % Channel 2 basin: x in [3.5, 5.5], floor at y = 1.2
    \draw[thick, fill=gray!20] (3.5, 1.2) rectangle (5.5, 0);
    \draw[thick] (3.5, 3.5) -- (3.5, 1.2) -- (5.5, 1.2) -- (5.5, 3.5);
    \node at (4.5, 0.6) {\textbf{Noise } $\sigma_2^2 = 0.6$};
    
    % Channel 3 basin: x in [6.5, 8.5], floor at y = 3.0
    \draw[thick, fill=gray!20] (6.5, 3.0) rectangle (8.5, 0);
    \draw[thick] (6.5, 3.5) -- (6.5, 3.0) -- (8.5, 3.0) -- (8.5, 3.5);
    \node at (7.5, 1.5) {\textbf{Noise } $\sigma_3^2 = 1.5$};

    % Water level mu* = 1.40 => y = 2.8
    \draw[very thick, dashed, color=coralborder] (-0.2, 2.8) -- (9.0, 2.8) 
        node[right, color=coralborder] {\textbf{Water Level } $\mu^* = \frac{1}{\nu^*} = 1.40$};

    % Water in Channel 1 (height 2.4, y from 0.4 to 2.8)
    \fill[tealborder!35] (0.52, 0.42) rectangle (2.48, 2.78);
    \node[color=primarydark, font=\bfseries] at (1.5, 1.6) {$P_1^* = 1.20$};

    % Water in Channel 2 (height 1.6, y from 1.2 to 2.8)
    \fill[tealborder!35] (3.52, 1.22) rectangle (5.48, 2.78);
    \node[color=primarydark, font=\bfseries] at (4.5, 2.0) {$P_2^* = 0.80$};

    % Channel 3 has floor at 3.0 > 2.8 => DRY!
    \node[color=traprad, font=\bfseries, align=center] at (7.5, 3.25) {\textsc{Dry Basin} \\ $P_3^* = 0.0$};

    % Labels
    \node[below=4pt] at (1.5, 0) {\textbf{Channel 1} (Clean)};
    \node[below=4pt] at (4.5, 0) {\textbf{Channel 2} (Moderate)};
    \node[below=4pt] at (7.5, 0) {\textbf{Channel 3} (Severe Noise)};
\end{tikzpicture}
\caption{The Water-Filling Principle for Parallel Gaussian Channels: Transmit power is poured into independent communication basins whose floor elevations represent noise power spectral densities $\sigma_i^2$. Channels with noise strictly below the equilibrium water level $\mu^*$ receive active transmit power $P_i^* = \mu^* - \sigma_i^2$, while channels with excessive noise ($\sigma_3^2 \ge \mu^*$) remain completely dry ($P_3^* = 0$), mathematically demonstrating how inequality constraints prune suboptimal dimensions.}
\label{fig:waterfilling_geometry}
\end{figure}

% ----------------------------------------------------------------------
% PART 4: DEEP INTUITION & CONCEPTUAL DERIVATIONS
% ----------------------------------------------------------------------
\section{Deep Intuition \& Conceptual Derivations}

\subsection{Projected Gradient Descent (PGD) and Non-Expansiveness}
When minimizing a differentiable convex loss functional $f: \mathbb{R}^d \to \mathbb{R}$ over a closed convex constraint set $C \subset \mathbb{R}^d$:
\begin{equation}
\min_{\mathbf{w} \in C} f(\mathbf{w})
\end{equation}
the \textbf{Projected Gradient Descent (PGD)} algorithm interleaves standard unconstrained gradient steps with orthogonal Euclidean projections back onto $C$:
\begin{equation}
\mathbf{w}_{k+1} = \Pi_C\left( \mathbf{w}_k - \eta \nabla f(\mathbf{w}_k) \right)
\label{eq:pgd_update}
\end{equation}

\begin{theorembox}{Variational Characterization of Euclidean Projection}
Let $C \subseteq \mathbb{R}^d$ be a non-empty closed convex set, and let $\mathbf{w} \in \mathbb{R}^d$. A point $\mathbf{z}^* \in C$ is the unique Euclidean projection of $\mathbf{w}$ onto $C$ ($\mathbf{z}^* = \Pi_C(\mathbf{w}) = \argmin_{\mathbf{z} \in C} \frac{1}{2}\|\mathbf{z} - \mathbf{w}\|_2^2$) if and only if:
\begin{equation}
\langle \mathbf{w} - \mathbf{z}^*, \mathbf{x} - \mathbf{z}^* \rangle \le 0 \quad \forall \mathbf{x} \in C
\label{eq:proj_variational_char}
\end{equation}
Geometrically, this obtuse angle condition ensures that moving from the projection $\mathbf{z}^*$ toward any other feasible point $\mathbf{x} \in C$ forms an angle $\ge 90^\circ$ with the displacement vector $\mathbf{w} - \mathbf{z}^*$.
\end{theorembox}

\begin{proof}
\textbf{Necessity ($\implies$):} Suppose $\mathbf{z}^* = \Pi_C(\mathbf{w})$. For any arbitrary point $\mathbf{x} \in C$, consider the convex combination $\mathbf{z}(t) = (1 - t)\mathbf{z}^* + t \mathbf{x} = \mathbf{z}^* + t(\mathbf{x} - \mathbf{z}^*)$ for $t \in [0, 1]$. Because $C$ is convex, $\mathbf{z}(t) \in C$ for all $t \in [0, 1]$. Define the scalar quadratic objective:
\begin{equation}
\phi(t) \triangleq \frac{1}{2}\|\mathbf{z}(t) - \mathbf{w}\|_2^2 = \frac{1}{2}\|\mathbf{z}^* + t(\mathbf{x} - \mathbf{z}^*) - \mathbf{w}\|_2^2
\end{equation}
Since $\mathbf{z}^* = \mathbf{z}(0)$ minimizes the distance over $C$, the one-sided directional derivative at $t = 0$ must be non-negative: $\phi'(0) \ge 0$. Differentiating with respect to $t$:
\begin{equation}
\phi'(t) = \langle \mathbf{z}(t) - \mathbf{w}, \mathbf{x} - \mathbf{z}^* \rangle \implies \phi'(0) = \langle \mathbf{z}^* - \mathbf{w}, \mathbf{x} - \mathbf{z}^* \rangle \ge 0
\end{equation}
Multiplying by $-1$ reverses the inequality:
\begin{equation}
\langle \mathbf{w} - \mathbf{z}^*, \mathbf{x} - \mathbf{z}^* \rangle \le 0
\end{equation}
\textbf{Sufficiency ($\impliedby$):} Suppose $\mathbf{z}^* \in C$ satisfies~\eqref{eq:proj_variational_char} for all $\mathbf{x} \in C$. For any feasible point $\mathbf{x} \in C$ with $\mathbf{x} \ne \mathbf{z}^*$, expand the squared distance:
\begin{align}
\|\mathbf{x} - \mathbf{w}\|_2^2 &= \|(\mathbf{x} - \mathbf{z}^*) + (\mathbf{z}^* - \mathbf{w})\|_2^2 \nonumber \\
&= \|\mathbf{x} - \mathbf{z}^*\|_2^2 + \|\mathbf{z}^* - \mathbf{w}\|_2^2 + 2\langle \mathbf{z}^* - \mathbf{w}, \mathbf{x} - \mathbf{z}^* \rangle \nonumber \\
&= \|\mathbf{x} - \mathbf{z}^*\|_2^2 + \|\mathbf{z}^* - \mathbf{w}\|_2^2 - 2\langle \mathbf{w} - \mathbf{z}^*, \mathbf{x} - \mathbf{z}^* \rangle
\end{align}
Because $-2\langle \mathbf{w} - \mathbf{z}^*, \mathbf{x} - \mathbf{z}^* \rangle \ge 0$ by~\eqref{eq:proj_variational_char} and $\|\mathbf{x} - \mathbf{z}^*\|_2^2 > 0$ for $\mathbf{x} \ne \mathbf{z}^*$:
\begin{equation}
\|\mathbf{x} - \mathbf{w}\|_2^2 > \|\mathbf{z}^* - \mathbf{w}\|_2^2
\end{equation}
Thus, $\mathbf{z}^*$ is the unique global minimizer of Euclidean distance over $C$.
\end{proof}

\begin{theorembox}{Firm Non-Expansiveness of the Convex Projection Operator}
Let $C \subseteq \mathbb{R}^d$ be a closed convex set. The Euclidean projection operator $\Pi_C: \mathbb{R}^d \to C$ is \textbf{firmly non-expansive}:
\begin{equation}
\|\Pi_C(\mathbf{u}) - \Pi_C(\mathbf{v})\|_2^2 + \|(I - \Pi_C)(\mathbf{u}) - (I - \Pi_C)(\mathbf{v})\|_2^2 \le \|\mathbf{u} - \mathbf{v}\|_2^2 \quad \forall \mathbf{u}, \mathbf{v} \in \mathbb{R}^d
\label{eq:firm_non_expansive}
\end{equation}
Consequently, $\Pi_C$ is strictly non-expansive ($\|\Pi_C(\mathbf{u}) - \Pi_C(\mathbf{v})\|_2 \le \|\mathbf{u} - \mathbf{v}\|_2$).
\end{theorembox}

\begin{proof}
Let $\hat{\mathbf{u}} = \Pi_C(\mathbf{u})$ and $\hat{\mathbf{v}} = \Pi_C(\mathbf{v})$. Applying the variational characterization~\eqref{eq:proj_variational_char}:
\begin{align}
\langle \mathbf{u} - \hat{\mathbf{u}}, \mathbf{x} - \hat{\mathbf{u}} \rangle &\le 0 \quad \forall \mathbf{x} \in C \label{eq:var_u} \\
\langle \mathbf{v} - \hat{\mathbf{v}}, \mathbf{x}' - \hat{\mathbf{v}} \rangle &\le 0 \quad \forall \mathbf{x}' \in C \label{eq:var_v}
\end{align}
Since $\hat{\mathbf{v}} \in C$, substitute $\mathbf{x} = \hat{\mathbf{v}}$ into~\eqref{eq:var_u}. Similarly, substitute $\mathbf{x}' = \hat{\mathbf{u}}$ into~\eqref{eq:var_v}:
\begin{align}
\langle \mathbf{u} - \hat{\mathbf{u}}, \hat{\mathbf{v}} - \hat{\mathbf{u}} \rangle &\le 0 \implies \langle \hat{\mathbf{u}} - \mathbf{u}, \hat{\mathbf{u}} - \hat{\mathbf{v}} \rangle \le 0 \\
\langle \mathbf{v} - \hat{\mathbf{v}}, \hat{\mathbf{u}} - \hat{\mathbf{v}} \rangle &\le 0
\end{align}
Summing these two inequalities:
\begin{equation}
\langle (\hat{\mathbf{u}} - \hat{\mathbf{v}}) - (\mathbf{u} - \mathbf{v}), \hat{\mathbf{u}} - \hat{\mathbf{v}} \rangle \le 0 \implies \|\hat{\mathbf{u}} - \hat{\mathbf{v}}\|_2^2 \le \langle \mathbf{u} - \mathbf{v}, \hat{\mathbf{u}} - \hat{\mathbf{v}} \rangle
\label{eq:firm_step}
\end{equation}
Now consider the vector identity for any vectors $\mathbf{a}, \mathbf{b} \in \mathbb{R}^d$: $\|\mathbf{a} + \mathbf{b}\|_2^2 = \|\mathbf{a}\|_2^2 + \|\mathbf{b}\|_2^2 + 2\langle \mathbf{a}, \mathbf{b} \rangle$. Setting $\mathbf{a} = \hat{\mathbf{u}} - \hat{\mathbf{v}}$ and $\mathbf{b} = (\mathbf{u} - \hat{\mathbf{u}}) - (\mathbf{v} - \hat{\mathbf{v}})$ so that $\mathbf{a} + \mathbf{b} = \mathbf{u} - \mathbf{v}$:
\begin{align}
\|\mathbf{u} - \mathbf{v}\|_2^2 &= \|\hat{\mathbf{u}} - \hat{\mathbf{v}}\|_2^2 + \|(\mathbf{u} - \hat{\mathbf{u}}) - (\mathbf{v} - \hat{\mathbf{v}})\|_2^2 + 2\langle \hat{\mathbf{u}} - \hat{\mathbf{v}}, (\mathbf{u} - \hat{\mathbf{u}}) - (\mathbf{v} - \hat{\mathbf{v}}) \rangle \nonumber \\
&= \|\Pi_C(\mathbf{u}) - \Pi_C(\mathbf{v})\|_2^2 + \|(I - \Pi_C)(\mathbf{u}) - (I - \Pi_C)(\mathbf{v})\|_2^2 + 2\big(\langle \hat{\mathbf{u}} - \hat{\mathbf{v}}, \mathbf{u} - \mathbf{v} \rangle - \|\hat{\mathbf{u}} - \hat{\mathbf{v}}\|_2^2\big)
\end{align}
By~\eqref{eq:firm_step}, the cross-term $\langle \hat{\mathbf{u}} - \hat{\mathbf{v}}, \mathbf{u} - \mathbf{v} \rangle - \|\hat{\mathbf{u}} - \hat{\mathbf{v}}\|_2^2 \ge 0$. Dropping this non-negative term yields~\eqref{eq:firm_non_expansive}. Dropping the non-negative residual norm $\|(I - \Pi_C)(\mathbf{u}) - (I - \Pi_C)(\mathbf{v})\|_2^2 \ge 0$ immediately gives $\|\Pi_C(\mathbf{u}) - \Pi_C(\mathbf{v})\|_2 \le \|\mathbf{u} - \mathbf{v}\|_2$.
\end{proof}

\begin{ponderbox}{Does Projection Accelerate or Retard Convergence?}
Since the projection operator satisfies $\|\Pi_C(\mathbf{u}) - \Pi_C(\mathbf{v})\|_2 \le \|\mathbf{u} - \mathbf{v}\|_2$, does the projection step shrink the distance to the constrained global minimizer $\mathbf{w}^* \in C$?
\tcblower
\textbf{The Socratic Resolution:} Yes! Notice that because the optimal solution $\mathbf{w}^*$ resides in $C$, projecting it leaves it completely unaltered: $\Pi_C(\mathbf{w}^*) = \mathbf{w}^*$. 

Setting $\mathbf{u} = \mathbf{w}_k - \eta \nabla f(\mathbf{w}_k)$ and $\mathbf{v} = \mathbf{w}^*$ in the non-expansiveness inequality:
\begin{align}
\|\mathbf{w}_{k+1} - \mathbf{w}^*\|_2 &= \|\Pi_C(\mathbf{w}_k - \eta \nabla f(\mathbf{w}_k)) - \Pi_C(\mathbf{w}^*)\|_2 \nonumber \\
&\le \|(\mathbf{w}_k - \eta \nabla f(\mathbf{w}_k)) - \mathbf{w}^*\|_2
\end{align}
The projection operator \textbf{never} pushes the parameter iterate farther away from the target $\mathbf{w}^*$ than the intermediate unconstrained step did. If the unconstrained gradient step overshoots the feasible domain, the orthogonal projection strictly contracts the distance to the optimum:
\begin{equation}
\|\mathbf{w}_{k+1} - \mathbf{w}^*\|_2^2 \le \|(\mathbf{w}_k - \eta \nabla f(\mathbf{w}_k)) - \mathbf{w}^*\|_2^2 - \|\mathbf{w}_{k+1} - (\mathbf{w}_k - \eta \nabla f(\mathbf{w}_k))\|_2^2
\end{equation}
This Pythagorean contraction ensures that Projected Gradient Descent achieves the identical $\mathcal{O}(1/k)$ sublinear convergence rate on convex sets as unconstrained Gradient Descent.
\end{ponderbox}

\subsection{Jensen's Inequality: The Engine of Probabilistic Learning}
For any convex functional $g: \mathbb{R} \to \mathbb{R}$ and any random variable $X$ defined on a probability space:
\begin{equation}
g(\E[X]) \le \E[g(X)]
\label{eq:jensen}
\end{equation}
Geometrically, Jensen's inequality is the statistical realization of secant chord convexity. If we take any probability mass distribution $p_1, \dots, p_n$ across values $x_1, \dots, x_n$, their center of mass $\sum p_i x_i = \E[X]$ evaluates to a point $g(\E[X])$ on the curve. In contrast, the expectation $\E[g(X)] = \sum p_i g(x_i)$ is the vertical center of mass of the corresponding points on the curve, which must float on or above the convex bowl.

In Machine Learning Foundations, Jensen's inequality serves as the primary theoretical engine for:
\begin{enumerate}
    \item \textbf{The Evidence Lower Bound (ELBO):} Derived in Chapter~\ref{chap:week12} for the Expectation-Maximization (EM) algorithm. Applying Jensen's inequality to the concave function $\log(\cdot)$:
    \begin{equation}
    \log p(\mathbf{x}) = \log \sum_{\mathbf{z}} q(\mathbf{z}) \frac{p(\mathbf{x}, \mathbf{z})}{q(\mathbf{z})} \ge \sum_{\mathbf{z}} q(\mathbf{z}) \log \frac{p(\mathbf{x}, \mathbf{z})}{q(\mathbf{z})} \triangleq \operatorname{ELBO}(q)
    \end{equation}
    \item \textbf{Non-Negativity of Kullback-Leibler Divergence:}
    \begin{equation}
    D_{\mathrm{KL}}(p \parallel q) = \E_p\left[ -\log \frac{q(X)}{p(X)} \right] \ge -\log \E_p\left[ \frac{q(X)}{p(X)} \right] = -\log(1) = 0
    \end{equation}
    guaranteeing that relative entropy is strictly non-negative, with equality if and only if $p = q$ almost everywhere.
\end{enumerate}

\begin{ponderbox}{The Water-Filling Paradox: Pruning Corrupted Channels}
In parallel Gaussian communication or multi-task resource allocation, common intuition suggests that every available channel should receive at least a small sliver of power ($\epsilon > 0$) to harvest its information throughput. Why does the optimal KKT solution ruthlessly assign $P_k^* = 0$ to noisy channels, and why would providing $\epsilon > 0$ power to those channels strictly degrade total transmission capacity?
\tcblower
\textbf{The Socratic Resolution:} This phenomenon is governed by the \textbf{Equal Marginal Returns Principle} of constrained convex optimization.

The marginal capacity gain obtained from injecting an incremental watt into channel $i$ is:
\begin{equation}
\frac{\partial C}{\partial P_i} = \frac{1}{\ln 2} \cdot \frac{1}{P_i + \sigma_i^2}
\end{equation}
Notice that when power is zero ($P_i = 0$), the initial marginal return is bounded above by the inverse noise floor: $\left.\frac{\partial C}{\partial P_i}\right|_{P_i = 0} = \frac{1}{\sigma_i^2 \ln 2}$. 

Now suppose the system reaches equilibrium across the clean channels at water level $\mu^* = 1/\nu^*$. Every active channel satisfies $P_j^* + \sigma_j^2 = \mu^*$, yielding identical marginal capacity returns: $\frac{\partial C}{\partial P_j} = \frac{1}{\mu^* \ln 2}$.

If a severely corrupted channel has noise variance $\sigma_k^2 > \mu^*$, its maximum possible marginal return (at $P_k = 0$) is:
\begin{equation}
\left.\frac{\partial C}{\partial P_k}\right|_{P_k = 0} = \frac{1}{\sigma_k^2 \ln 2} < \frac{1}{\mu^* \ln 2}
\end{equation}
Because the total power budget $\sum P_i = P_{\text{total}}$ is strictly conserved, granting $\epsilon > 0$ power to the noisy channel requires stealing $\epsilon$ power from clean channels. 
The capacity loss incurred by depriving clean channels is $\approx \frac{\epsilon}{\mu^* \ln 2}$, whereas the capacity gained from the noisy channel is only $\approx \frac{\epsilon}{\sigma_k^2 \ln 2}$. Because $\sigma_k^2 > \mu^*$:
\begin{equation}
\Delta C_{\text{total}} \approx \epsilon \left( \frac{1}{\sigma_k^2 \ln 2} - \frac{1}{\mu^* \ln 2} \right) < 0
\end{equation}
Every watt diverted to the noisy channel produces less capacity than it destroys in the clean channels. In constrained optimization, the boundary condition $P_k^* = 0$ is not a defect; it is nature's mechanism for pruning suboptimal dimensions.
\end{ponderbox}

% ----------------------------------------------------------------------
% PART 5: MULTI-PAGE WORKED EXAMPLES
% ----------------------------------------------------------------------
\section{Worked Numerical Examples}

\begin{examplebox}[breakable]{Step-by-Step Lagrange Multiplier Optimization on an Anisotropic Objective}
\textbf{Problem Statement:}
Consider minimizing the anisotropic quadratic objective functional:
\begin{equation}
f(x, y) = 2x^2 + y^2
\end{equation}
subject to the linear equality constraint:
\begin{equation}
h(x, y) = x + y - 3 = 0
\end{equation}

\textbf{Task Requirements:}
\begin{enumerate}
    \item Formulate the Lagrangian functional $\mathcal{L}(x, y, \nu)$.
    \item Compute the partial derivatives and set up the stationarity linear system.
    \item Solve for the exact optimal coordinates $(x^*, y^*)$ and the optimal multiplier $\nu^*$.
    \item Verify that the objective gradient $\nabla f(x^*, y^*)$ and constraint gradient $\nabla h(x^*, y^*)$ satisfy the collinearity condition.
    \item Compute the optimal objective value $f(x^*, y^*)$ and rigorously contrast it with the value obtained by projecting the unconstrained minimizer $(0, 0)$ onto $x + y = 3$.
\end{enumerate}

\tcblower

\textbf{Step 1: Formulation of the Lagrangian}
The Lagrangian introduces multiplier $\nu \in \mathbb{R}$:
\begin{equation}
\mathcal{L}(x, y, \nu) = f(x, y) + \nu \, h(x, y) = 2x^2 + y^2 + \nu(x + y - 3)
\end{equation}

\textbf{Step 2: Stationarity Conditions}
Setting all partial derivatives to zero:
\begin{align}
\frac{\partial \mathcal{L}}{\partial x} &= 4x + \nu = 0 \label{eq:ex_stat_x} \\
\frac{\partial \mathcal{L}}{\partial y} &= 2y + \nu = 0 \label{eq:ex_stat_y} \\
\frac{\partial \mathcal{L}}{\partial \nu} &= x + y - 3 = 0 \label{eq:ex_stat_nu}
\end{align}

\textbf{Step 3: Analytical Solution}
From~\eqref{eq:ex_stat_x} and~\eqref{eq:ex_stat_y}, express $x$ and $y$ as functions of $\nu$:
\begin{align}
x &= -\frac{\nu}{4} \\
y &= -\frac{\nu}{2}
\end{align}
Substitute these expressions into the primal feasibility equation~\eqref{eq:ex_stat_nu}:
\begin{equation}
-\frac{\nu}{4} - \frac{\nu}{2} - 3 = 0 \implies -\frac{3\nu}{4} = 3 \implies \nu^* = -4
\end{equation}
Now back-substitute $\nu^* = -4$:
\begin{align}
x^* &= -\frac{-4}{4} = 1.0 \\
y^* &= -\frac{-4}{2} = 2.0
\end{align}
The optimal constrained solution is $\mathbf{w}^* = (x^*, y^*) = (1, 2)$.

\textbf{Step 4: Verification of Gradient Collinearity}
Compute the gradient of the objective function at $(1, 2)$:
\begin{equation}
\nabla f(1, 2) = \begin{bmatrix} 4x \\ 2y \end{bmatrix}\Bigg|_{(1, 2)} = \begin{bmatrix} 4(1) \\ 2(2) \end{bmatrix} = \begin{bmatrix} 4 \\ 4 \end{bmatrix}
\end{equation}
Compute the gradient of the constraint function:
\begin{equation}
\nabla h(1, 2) = \begin{bmatrix} 1 \\ 1 \end{bmatrix}
\end{equation}
Check the stationarity condition:
\begin{equation}
\nabla f(1, 2) + \nu^* \nabla h(1, 2) = \begin{bmatrix} 4 \\ 4 \end{bmatrix} + (-4) \begin{bmatrix} 1 \\ 1 \end{bmatrix} = \begin{bmatrix} 0 \\ 0 \end{bmatrix}
\end{equation}
Equivalently:
\begin{equation}
\nabla f(1, 2) = 4 \begin{bmatrix} 1 \\ 1 \end{bmatrix} = -\nu^* \nabla h(1, 2)
\end{equation}
The gradients are strictly collinear. The objective contour ellipse $2x^2 + y^2 = 6$ is exactly tangent to the constraint line $x + y = 3$ at $(1, 2)$.

\textbf{Step 5: Optimal Value vs.\ Unconstrained Projection}
The true constrained optimal value is:
\begin{equation}
f(x^*, y^*) = 2(1)^2 + (2)^2 = 2 + 4 = 6.000
\end{equation}
Now consider the naive projection heuristic:
The unconstrained minimizer is $(x_{\text{unc}}, y_{\text{unc}}) = (0, 0)$.
The Euclidean projection of $(0, 0)$ onto the hyperplane $x + y = 3$ is:
\begin{equation}
\Pi(0, 0) = \mathbf{0} - \frac{\mathbf{a}^T \mathbf{0} - b}{\|\mathbf{a}\|_2^2} \mathbf{a} = \mathbf{0} - \frac{0 + 0 - 3}{1^2 + 1^2} \begin{bmatrix} 1 \\ 1 \end{bmatrix} = \frac{3}{2} \begin{bmatrix} 1 \\ 1 \end{bmatrix} = (1.5, 1.5)
\end{equation}
Evaluating the objective at this projected point:
\begin{equation}
f(1.5, 1.5) = 2(1.5)^2 + (1.5)^2 = 2(2.25) + 2.25 = 4.5 + 2.25 = 6.750
\end{equation}
Suboptimality gap:
\begin{equation}
\Delta f = 6.750 - 6.000 = 0.750 \quad (+12.5\% \text{ excess loss})
\end{equation}
This rigorous calculation confirms that projecting the unconstrained minimum onto the feasible set fails because Euclidean projection treats all coordinate axes symmetrically, whereas the objective function is twice as steep along the $x$-axis as along the $y$-axis.
\end{examplebox}

\begin{examplebox}[breakable]{Convexity Analysis and Hessian Nullspace of the Log-Sum-Exp Function}
\textbf{Problem Statement:}
In multinomial logistic regression and deep neural networks, the cross-entropy loss relies on the log-sum-exp functional:
\begin{equation}
f(\mathbf{x}) = \log\left( \sum_{i=1}^d \exp(x_i) \right) \quad \text{for } \mathbf{x} \in \mathbb{R}^d
\end{equation}

\textbf{Task Requirements:}
\begin{enumerate}
    \item Compute the gradient $\nabla f(\mathbf{x})$ and prove it forms the softmax probability vector.
    \item Derive the $d \times d$ Hessian matrix $\nabla^2 f(\mathbf{x})$.
    \item Prove that $\nabla^2 f(\mathbf{x}) \succeq 0$ for all $\mathbf{x} \in \mathbb{R}^d$ using statistical variance identities.
    \item Characterize the nullspace of $\nabla^2 f(\mathbf{x})$ and explain why $f(\mathbf{x})$ is convex, but \emph{not strictly convex}.
\end{enumerate}

\tcblower

\textbf{Step 1: Gradient Derivation (The Softmax Mapping)}
Let $S(\mathbf{x}) = \sum_{j=1}^d e^{x_j}$. Then $f(\mathbf{x}) = \log S(\mathbf{x})$.
By the multivariable chain rule, for any coordinate $k \in \{1, \dots, d\}$:
\begin{equation}
\frac{\partial f}{\partial x_k} = \frac{1}{S(\mathbf{x})} \frac{\partial S(\mathbf{x})}{\partial x_k} = \frac{e^{x_k}}{\sum_{j=1}^d e^{x_j}} \triangleq \sigma_k(\mathbf{x})
\end{equation}
Observe that:
\begin{equation}
\sigma_k(\mathbf{x}) > 0 \quad \forall k, \quad \text{and} \quad \sum_{k=1}^d \sigma_k(\mathbf{x}) = \frac{\sum_{k=1}^d e^{x_k}}{S(\mathbf{x})} = 1
\end{equation}
Thus, the gradient vector $\nabla f(\mathbf{x}) = \boldsymbol{\sigma}(\mathbf{x}) = [\sigma_1, \dots, \sigma_d]^T$ lies strictly within the interior of the probability simplex $\Delta^d$.

\textbf{Step 2: Hessian Matrix Derivation}
Differentiating $\sigma_k$ with respect to $x_j$:
\begin{itemize}
    \item When $j = k$:
    \begin{equation}
    \frac{\partial \sigma_k}{\partial x_k} = \frac{e^{x_k} S - e^{x_k} e^{x_k}}{S^2} = \frac{e^{x_k}}{S} - \left(\frac{e^{x_k}}{S}\right)^2 = \sigma_k - \sigma_k^2 = \sigma_k(1 - \sigma_k)
    \end{equation}
    \item When $j \ne k$:
    \begin{equation}
    \frac{\partial \sigma_k}{\partial x_j} = \frac{0 \cdot S - e^{x_k} e^{x_j}}{S^2} = -\sigma_k \sigma_j
    \end{equation}
\end{itemize}
Using the Kronecker delta $\delta_{kj}$:
\begin{equation}
[\nabla^2 f(\mathbf{x})]_{kj} = \sigma_k \delta_{kj} - \sigma_k \sigma_j
\end{equation}
In compact matrix form:
\begin{equation}
\nabla^2 f(\mathbf{x}) = \diag(\boldsymbol{\sigma}) - \boldsymbol{\sigma}\boldsymbol{\sigma}^T
\end{equation}

\textbf{Step 3: Proof of Positive Semidefiniteness via Variance}
Let $\mathbf{v} \in \mathbb{R}^d$ be an arbitrary non-zero test vector. Compute the quadratic form:
\begin{align}
\mathbf{v}^T \nabla^2 f(\mathbf{x}) \mathbf{v} &= \mathbf{v}^T \left( \diag(\boldsymbol{\sigma}) - \boldsymbol{\sigma}\boldsymbol{\sigma}^T \right) \mathbf{v} \nonumber \\
&= \mathbf{v}^T \diag(\boldsymbol{\sigma}) \mathbf{v} - (\mathbf{v}^T \boldsymbol{\sigma})^2 \nonumber \\
&= \sum_{i=1}^d \sigma_i v_i^2 - \left( \sum_{i=1}^d \sigma_i v_i \right)^2
\end{align}
Interpret the components of $\mathbf{v}$ as values taken by a discrete random variable $V$, with probability distribution given by $\Pr(V = v_i) = \sigma_i$.
Then:
\begin{align}
\sum_{i=1}^d \sigma_i v_i^2 &= \E[V^2] \\
\sum_{i=1}^d \sigma_i v_i &= \E[V]
\end{align}
Substituting these statistical expectations into the quadratic form:
\begin{equation}
\mathbf{v}^T \nabla^2 f(\mathbf{x}) \mathbf{v} = \E[V^2] - (\E[V])^2 = \Var(V)
\end{equation}
Because the mathematical variance of any real random variable is non-negative ($\Var(V) \ge 0$):
\begin{equation}
\mathbf{v}^T \nabla^2 f(\mathbf{x}) \mathbf{v} \ge 0 \quad \forall \mathbf{v} \in \mathbb{R}^d \implies \nabla^2 f(\mathbf{x}) \succeq 0
\end{equation}
This rigorously establishes that the log-sum-exp functional is globally convex on $\mathbb{R}^d$.

\textbf{Step 4: Characterization of the Nullspace}
When does $\Var(V) = 0$?
The variance of a discrete random variable is zero if and only if $V$ is a constant almost surely:
\begin{equation}
v_1 = v_2 = \dots = v_d = c \in \mathbb{R} \iff \mathbf{v} = c \mathbf{1}
\end{equation}
where $\mathbf{1} = [1, 1, \dots, 1]^T$.
Therefore, the nullspace of the Hessian is one-dimensional:
\begin{equation}
\Null(\nabla^2 f(\mathbf{x})) = \Span(\mathbf{1})
\end{equation}
For any shift direction $\mathbf{v} = c\mathbf{1}$:
\begin{equation}
f(\mathbf{x} + c\mathbf{1}) = \log\left(\sum_{i=1}^d e^{x_i + c}\right) = \log\left(e^c \sum_{i=1}^d e^{x_i}\right) = c + \log\left(\sum_{i=1}^d e^{x_i}\right) = f(\mathbf{x}) + c
\end{equation}
Along the line of ones, $f$ is purely linear (affine), which confirms that $\nabla^2 f$ has an eigenvalue of zero: $\nabla^2 f(\mathbf{x}) \mathbf{1} = \mathbf{0}$.
Consequently, $f$ is \textbf{convex}, but \textbf{not strictly convex}. It is strictly convex only on the orthogonal subspace $\{\mathbf{v} \mid \mathbf{1}^T \mathbf{v} = 0\}$.
\end{examplebox}

\begin{examplebox}[breakable]{Exact Euclidean Projection onto a Polyhedral Halfspace}
\textbf{Problem Statement:}
Consider the closed polyhedral halfspace:
\begin{equation}
C = \left\{ \mathbf{x} \in \mathbb{R}^2 \;\middle|\; 2x_1 + x_2 \le 4 \right\}
\end{equation}
Let $\mathbf{x}_0 = [3.0, 4.0]^T \in \mathbb{R}^2$ be an external test point.

\textbf{Task Requirements:}
\begin{enumerate}
    \item Verify that $\mathbf{x}_0$ is infeasible and violates the halfspace constraint.
    \item Formulate the projection problem as a constrained quadratic minimization task.
    \item Solve for the exact projection coordinates $\mathbf{x}^* = \Pi_C(\mathbf{x}_0)$ using Lagrange multipliers.
    \item Verify the orthogonality condition: show that the displacement vector $\mathbf{x}_0 - \mathbf{x}^*$ is parallel to the constraint normal vector $\mathbf{a} = [2, 1]^T$.
    \item Calculate the exact Euclidean distance from $\mathbf{x}_0$ to the halfspace $C$.
\end{enumerate}

\tcblower

\textbf{Step 1: Feasibility Verification}
Evaluate the constraint functional at $\mathbf{x}_0 = [3.0, 4.0]^T$:
\begin{equation}
g(\mathbf{x}_0) = 2(3.0) + (4.0) = 6.0 + 4.0 = 10.0 > 4.0
\end{equation}
Because $10.0 > 4.0$, $\mathbf{x}_0 \notin C$. Since $C$ is a closed convex set, the projection $\mathbf{x}^*$ must lie on the active boundary hyperplane:
\begin{equation}
\mathcal{H} = \{\mathbf{x} \in \mathbb{R}^2 \mid 2x_1 + x_2 = 4\}
\end{equation}

\textbf{Step 2: Constrained Quadratic Minimization Formulation}
The Euclidean projection solves:
\begin{equation}
\mathbf{x}^* = \argmin_{\mathbf{x} \in \mathbb{R}^2} \frac{1}{2} \|\mathbf{x} - \mathbf{x}_0\|_2^2 \quad \text{subject to } \mathbf{a}^T \mathbf{x} = b
\end{equation}
where $\mathbf{a} = [2, 1]^T$ and $b = 4$.
The Lagrangian functional is:
\begin{equation}
\mathcal{L}(\mathbf{x}, \nu) = \frac{1}{2}\|\mathbf{x} - \mathbf{x}_0\|_2^2 + \nu(\mathbf{a}^T \mathbf{x} - b)
\end{equation}

\textbf{Step 3: Stationarity and Analytical Solution}
Differentiating with respect to $\mathbf{x}$:
\begin{equation}
\nabla_{\mathbf{x}}\mathcal{L} = \mathbf{x} - \mathbf{x}_0 + \nu \mathbf{a} = \mathbf{0} \implies \mathbf{x} = \mathbf{x}_0 - \nu \mathbf{a}
\end{equation}
Multiply by $\mathbf{a}^T$ from the left to enforce the boundary constraint $\mathbf{a}^T \mathbf{x} = b$:
\begin{equation}
\mathbf{a}^T (\mathbf{x}_0 - \nu \mathbf{a}) = b \implies \mathbf{a}^T \mathbf{x}_0 - \nu \|\mathbf{a}\|_2^2 = b
\end{equation}
Solving for the optimal multiplier $\nu^*$:
\begin{equation}
\nu^* = \frac{\mathbf{a}^T \mathbf{x}_0 - b}{\|\mathbf{a}\|_2^2}
\end{equation}
Numerically evaluating the components:
\begin{align}
\mathbf{a}^T \mathbf{x}_0 - b &= 2(3.0) + 1(4.0) - 4 = 10.0 - 4 = 6.0 \\
\|\mathbf{a}\|_2^2 &= 2^2 + 1^2 = 4 + 1 = 5.0
\end{align}
Thus:
\begin{equation}
\nu^* = \frac{6.0}{5.0} = 1.2
\end{equation}
Now compute the projection vector $\mathbf{x}^*$:
\begin{align}
\mathbf{x}^* &= \mathbf{x}_0 - \nu^* \mathbf{a} \nonumber \\
&= \begin{bmatrix} 3.0 \\ 4.0 \end{bmatrix} - 1.2 \begin{bmatrix} 2.0 \\ 1.0 \end{bmatrix}
= \begin{bmatrix} 3.0 - 2.4 \\ 4.0 - 1.2 \end{bmatrix} = \begin{bmatrix} 0.6 \\ 2.8 \end{bmatrix}
\end{align}
The projected coordinates are $\mathbf{x}^* = (0.6, 2.8)$.

\textbf{Step 4: Boundary and Orthogonality Verification}
\begin{itemize}
    \item Boundary check: $2(0.6) + 2.8 = 1.2 + 2.8 = 4.0$. Feasibility holds with exact equality.
    \item Orthogonality check: The displacement vector is:
    \begin{equation}
    \mathbf{d} = \mathbf{x}_0 - \mathbf{x}^* = \begin{bmatrix} 3.0 - 0.6 \\ 4.0 - 2.8 \end{bmatrix} = \begin{bmatrix} 2.4 \\ 1.2 \end{bmatrix} = 1.2 \begin{bmatrix} 2 \\ 1 \end{bmatrix} = 1.2 \, \mathbf{a}
    \end{equation}
    The displacement vector $\mathbf{d}$ is strictly parallel to the outward normal $\mathbf{a}$.
    Any tangent direction along the line $2x_1 + x_2 = 4$ is of the form $\mathbf{t} = [-1, 2]^T$. Testing orthogonality:
    \begin{equation}
    \mathbf{d}^T \mathbf{t} = \begin{bmatrix} 2.4 & 1.2 \end{bmatrix} \begin{bmatrix} -1 \\ 2 \end{bmatrix} = -2.4 + 2.4 = 0
    \end{equation}
\end{itemize}

\textbf{Step 5: Exact Euclidean Distance}
The minimum distance from $\mathbf{x}_0$ to the halfspace $C$ is:
\begin{equation}
\operatorname{dist}(\mathbf{x}_0, C) = \|\mathbf{x}_0 - \mathbf{x}^*\|_2 = \|\mathbf{d}\|_2 = \sqrt{(2.4)^2 + (1.2)^2} = \sqrt{5.76 + 1.44} = \sqrt{7.20} \approx 2.6833
\end{equation}
Notice that this matches the classical point-to-plane distance formula:
\begin{equation}
\operatorname{dist}(\mathbf{x}_0, \mathcal{H}) = \frac{|\mathbf{a}^T \mathbf{x}_0 - b|}{\|\mathbf{a}\|_2} = \frac{|10 - 4|}{\sqrt{5}} = \frac{6}{\sqrt{5}} = \frac{6\sqrt{5}}{5} \approx 2.6833
\end{equation}
\end{examplebox}

\begin{examplebox}[breakable]{The Water-Filling Algorithm for Parallel Gaussian Channels}
\textbf{Problem Statement:}
In wireless communications and optimal resource allocation, a transmitter must distribute a total power budget $P_{\text{total}} = 2.0$ watts across $K = 3$ independent parallel Gaussian channels to maximize the total Shannon information throughput. The channel noise variances are:
\begin{equation}
\sigma_1^2 = 0.20, \quad \sigma_2^2 = 0.60, \quad \sigma_3^2 = 1.50
\end{equation}
The total Shannon capacity (in bits per second per Hertz) is given by:
\begin{equation}
C(P_1, P_2, P_3) = \frac{1}{2} \sum_{i=1}^3 \log_2\left( 1 + \frac{P_i}{\sigma_i^2} \right)
\end{equation}
Subject to the budget and non-negativity constraints:
\begin{equation}
\sum_{i=1}^3 P_i = 2.0, \quad P_i \ge 0 \quad (i = 1, 2, 3)
\end{equation}

\textbf{Task Requirements:}
\begin{enumerate}
    \item Formulate the equivalent convex minimization problem and its Lagrangian functional $\mathcal{L}(\mathbf{P}, \nu, \boldsymbol{\lambda})$.
    \item Derive the Karush-Kuhn-Tucker (KKT) first-order optimality conditions and the explicit water-filling equation $P_i^* = \max(0, \mu^* - \sigma_i^2)$ with water level $\mu^* \triangleq 1/\nu^*$.
    \item Solve step-by-step for the exact equilibrium water level $\mu^*$, proving why Channel~3 must be deactivated ($P_3^* = 0$).
    \item Compute the exact optimal power allocation $\mathbf{P}^* = [P_1^*, P_2^*, P_3^*]^T$ and dual multipliers $\nu^*$ and $\lambda_3^*$.
    \item Calculate the optimal Shannon capacity $C(\mathbf{P}^*)$ and contrast it with a naive uniform allocation $P_1 = P_2 = P_3 = \frac{2}{3}$, quantifying the percentage spectral efficiency gain.
\end{enumerate}

\tcblower

\textbf{Step 1: Convex Minimization Formulation \& Lagrangian}
Maximizing $\sum \frac{1}{2}\log_2(1 + P_i/\sigma_i^2)$ is mathematically equivalent to minimizing the negative natural logarithm:
\begin{equation}
\min_{\mathbf{P} \in \mathbb{R}^3} f(\mathbf{P}) \triangleq -\sum_{i=1}^3 \ln(P_i + \sigma_i^2) \quad \text{subject to } \sum_{i=1}^3 P_i = 2.0, \quad -P_i \le 0 \; (i=1, 2, 3)
\end{equation}
where we dropped the constant terms $\ln(\sigma_i^2)$. The objective is strictly convex because each $-\ln(P_i + \sigma_i^2)$ has positive second derivative $\frac{1}{(P_i + \sigma_i^2)^2} > 0$.
Introducing equality multiplier $\nu \in \mathbb{R}$ and inequality multipliers $\lambda_i \ge 0$:
\begin{equation}
\mathcal{L}(\mathbf{P}, \nu, \boldsymbol{\lambda}) = -\sum_{i=1}^3 \ln(P_i + \sigma_i^2) + \nu \left( \sum_{i=1}^3 P_i - 2.0 \right) - \sum_{i=1}^3 \lambda_i P_i
\end{equation}

\textbf{Step 2: KKT Optimality Conditions}
The KKT conditions require:
\begin{align}
\text{Stationarity: } & \frac{\partial \mathcal{L}}{\partial P_i} = -\frac{1}{P_i + \sigma_i^2} + \nu - \lambda_i = 0 \implies \frac{1}{P_i + \sigma_i^2} = \nu - \lambda_i \label{eq:wf_stat} \\
\text{Primal Feasibility: } & \sum_{i=1}^3 P_i = 2.0, \quad P_i \ge 0 \quad \forall i \in \{1, 2, 3\} \\
\text{Dual Feasibility: } & \lambda_i \ge 0 \quad \forall i \in \{1, 2, 3\} \\
\text{Complementary Slackness: } & \lambda_i P_i = 0 \quad \forall i \in \{1, 2, 3\} \label{eq:wf_slack}
\end{align}
From complementary slackness~\eqref{eq:wf_slack}:
\begin{itemize}
    \item If $P_i > 0$, then $\lambda_i = 0$. From~\eqref{eq:wf_stat}: $\frac{1}{P_i + \sigma_i^2} = \nu \implies P_i + \sigma_i^2 = \frac{1}{\nu} \triangleq \mu^*$. Thus $P_i = \mu^* - \sigma_i^2$.
    \item If $P_i = 0$, then $\lambda_i = \nu - \frac{1}{\sigma_i^2} \ge 0 \implies \nu \ge \frac{1}{\sigma_i^2} \iff \sigma_i^2 \ge \frac{1}{\nu} = \mu^*$.
\end{itemize}
Combining both cases yields the master \textbf{Water-Filling Formula}:
\begin{equation}
P_i^* = \max\left(0, \, \mu^* - \sigma_i^2\right) = (\mu^* - \sigma_i^2)^+
\end{equation}

\textbf{Step 3: Analytical Determination of Water Level $\mu^*$}
We determine $\mu^*$ by sorting noise levels in ascending order: $\sigma_1^2 = 0.20 < \sigma_2^2 = 0.60 < \sigma_3^2 = 1.50$.

\emph{Trial 1: Assume all $K=3$ channels are active ($P_i > 0$ for $i=1, 2, 3$).}
The total power constraint becomes:
\begin{equation}
\sum_{i=1}^3 (\mu - \sigma_i^2) = 3\mu - (0.20 + 0.60 + 1.50) = 3\mu - 2.30 = 2.00
\end{equation}
Solving for $\mu$:
\begin{equation}
3\mu = 4.30 \implies \mu = \frac{4.30}{3} \approx 1.4333
\end{equation}
Now check feasibility for Channel~3:
\begin{equation}
P_3 = \mu - \sigma_3^2 = 1.4333 - 1.50 = -0.0667 < 0
\end{equation}
This violates the primal non-negativity constraint $P_3 \ge 0$. Hence, Channel~3 cannot be active.

\emph{Trial 2: Set $P_3^* = 0$ (Channel~3 deactivated). Assume Channels~1 and~2 are active.}
The power budget must be satisfied exclusively by Channels~1 and~2:
\begin{equation}
(\mu - \sigma_1^2) + (\mu - \sigma_2^2) = (\mu - 0.20) + (\mu - 0.60) = 2\mu - 0.80 = 2.00
\end{equation}
Solving for $\mu^*$:
\begin{equation}
2\mu^* = 2.80 \implies \mu^* = 1.40
\end{equation}
Now verify consistency across all three channels:
\begin{align}
P_1^* &= \mu^* - \sigma_1^2 = 1.40 - 0.20 = 1.20 > 0 \quad \checkmark \\
P_2^* &= \mu^* - \sigma_2^2 = 1.40 - 0.60 = 0.80 > 0 \quad \checkmark \\
\sigma_3^2 &= 1.50 \ge \mu^* = 1.40 \implies P_3^* = 0.00 \quad \checkmark
\end{align}
Total power allocated: $P_1^* + P_2^* + P_3^* = 1.20 + 0.80 + 0.00 = 2.00 = P_{\text{total}}$ ($\checkmark$).

\textbf{Step 4: Verification of Dual Multipliers}
\begin{align}
\text{Equality Multiplier: } & \nu^* = \frac{1}{\mu^*} = \frac{1}{1.40} = \frac{5}{7} \approx 0.7143 \\
\text{Channel 1 Multiplier: } & \lambda_1^* = 0.00 \quad (\text{since } P_1^* > 0) \\
\text{Channel 2 Multiplier: } & \lambda_2^* = 0.00 \quad (\text{since } P_2^* > 0) \\
\text{Channel 3 Multiplier: } & \lambda_3^* = \nu^* - \frac{1}{\sigma_3^2} = \frac{5}{7} - \frac{1}{1.50} = \frac{5}{7} - \frac{2}{3} = \frac{15 - 14}{21} = \frac{1}{21} \approx 0.0476 \ge 0 \quad \checkmark
\end{align}
Because $\lambda_3^* > 0$, dual feasibility holds strictly, proving that Channel~3 is correctly inactive.

\textbf{Step 5: Capacity Evaluation \& Benchmark Comparison}
Compute the optimal Shannon capacity under water-filling allocation $\mathbf{P}^* = [1.20, 0.80, 0.00]^T$:
\begin{align}
C(\mathbf{P}^*) &= \frac{1}{2}\left[ \log_2\left(1 + \frac{1.20}{0.20}\right) + \log_2\left(1 + \frac{0.80}{0.60}\right) + \log_2\left(1 + \frac{0.00}{1.50}\right) \right] \nonumber \\
&= \frac{1}{2}\left[ \log_2(7.0) + \log_2(2.3333) + 0 \right] \nonumber \\
&= \frac{1}{2}\left[ 2.8074 + 1.2224 \right] = \frac{4.0298}{2} = 2.0149 \text{ bits/s/Hz}
\end{align}

Now contrast with a naive uniform allocation dividing $P_{\text{total}} = 2.0$ equally: $P_1 = P_2 = P_3 = \frac{2}{3} \approx 0.6667$:
\begin{align}
C(\mathbf{P}_{\text{uniform}}) &= \frac{1}{2}\left[ \log_2\left(1 + \frac{2/3}{0.20}\right) + \log_2\left(1 + \frac{2/3}{0.60}\right) + \log_2\left(1 + \frac{2/3}{1.50}\right) \right] \nonumber \\
&= \frac{1}{2}\left[ \log_2(1 + 3.3333) + \log_2(1 + 1.1111) + \log_2(1 + 0.4444) \right] \nonumber \\
&= \frac{1}{2}\left[ \log_2(4.3333) + \log_2(2.1111) + \log_2(1.4444) \right] \nonumber \\
&= \frac{1}{2}\left[ 2.1155 + 1.0780 + 0.5305 \right] = \frac{3.7240}{2} = 1.8620 \text{ bits/s/Hz}
\end{align}
The relative capacity gain delivered by optimal water-filling is:
\begin{equation}
\Delta C = \frac{2.0149 - 1.8620}{1.8620} \times 100\% = +8.21\%
\end{equation}
By recognizing that Channel~3's noise floor ($\sigma_3^2 = 1.50$) exceeds the system water level ($\mu^* = 1.40$) and shutting it down entirely, the transmitter diverts critical energy to Channels~1 and~2, achieving a substantial $8.21\%$ increase in spectral efficiency.
\end{examplebox}

\begin{examplebox}[breakable]{Step-by-Step Numerical Trace of Projected Gradient Descent}
\textbf{Problem Statement:}
Consider minimizing the strongly convex, anisotropic quadratic loss functional:
\begin{equation}
f(w_1, w_2) = (w_1 - 3)^2 + 4(w_2 - 2)^2
\end{equation}
subject to the hypercube box constraint set:
\begin{equation}
\mathcal{C} = \left\{ \mathbf{w} \in \mathbb{R}^2 \;\middle|\; 0 \le w_1 \le 1, \; 0 \le w_2 \le 1 \right\} = [0, 1] \times [0, 1]
\end{equation}
Notice that the unconstrained global minimizer is at $\mathbf{w}_{\text{unc}}^* = (3, 2)$ where $f(3, 2) = 0$, which lies strictly outside the feasible box $\mathcal{C}$.

\textbf{Task Requirements:}
\begin{enumerate}
    \item Compute the objective gradient $\nabla f(\mathbf{w})$ and state the coordinate-wise Euclidean projection operator $\Pi_{\mathcal{C}}(\mathbf{z})$.
    \item Perform the complete numerical calculation for **Iteration 1** starting from $\mathbf{w}_0 = (0.0, 0.0)^T$ with learning rate $\eta = 0.20$: compute the gradient $\nabla f(\mathbf{w}_0)$, unconstrained update $\tilde{\mathbf{w}}_1$, projected point $\mathbf{w}_1$, and loss $f(\mathbf{w}_1)$.
    \item Perform the complete numerical calculation for **Iteration 2**: compute $\nabla f(\mathbf{w}_1)$, unconstrained candidate $\tilde{\mathbf{w}}_2$, projected point $\mathbf{w}_2$, and test for numerical convergence.
    \item Rigorously prove that the terminal iterate satisfies the Variational Characterization of Optimality for constrained convex minimization:
    \begin{equation}
    \nabla f(\mathbf{w}^*)^T (\mathbf{z} - \mathbf{w}^*) \ge 0 \quad \forall \mathbf{z} \in \mathcal{C}
    \end{equation}
    confirming that the iterate is the exact, global constrained minimizer.
\end{enumerate}

\tcblower

\textbf{Step 1: Gradient Vector and Projection Operator}
The gradient of $f$ is:
\begin{equation}
\nabla f(\mathbf{w}) = \begin{bmatrix} \frac{\partial f}{\partial w_1} \\ \frac{\partial f}{\partial w_2} \end{bmatrix} = \begin{bmatrix} 2(w_1 - 3) \\ 8(w_2 - 2) \end{bmatrix}
\end{equation}
The Euclidean projection onto the box $\mathcal{C} = [0, 1] \times [0, 1]$ decouples into independent 1D interval projections (clipping):
\begin{equation}
\Pi_{\mathcal{C}}(\mathbf{z}) = \begin{bmatrix} \operatorname{clip}(z_1, 0, 1) \\ \operatorname{clip}(z_2, 0, 1) \end{bmatrix} = \begin{bmatrix} \min(1, \max(0, z_1)) \\ \min(1, \max(0, z_2)) \end{bmatrix}
\end{equation}

\textbf{Step 2: Iteration 1 Walkthrough}
\begin{itemize}
    \item Initial state: $\mathbf{w}_0 = [0.0, 0.0]^T$.
    \item Initial objective value:
    \begin{equation}
    f(\mathbf{w}_0) = (0 - 3)^2 + 4(0 - 2)^2 = (-3)^2 + 4(-2)^2 = 9 + 4(4) = 9 + 16 = 25.0
    \end{equation}
    \item Gradient computation at $\mathbf{w}_0$:
    \begin{equation}
    \nabla f(\mathbf{w}_0) = \begin{bmatrix} 2(0 - 3) \\ 8(0 - 2) \end{bmatrix} = \begin{bmatrix} -6.0 \\ -16.0 \end{bmatrix}
    \end{equation}
    \item Intermediate unconstrained gradient step ($\eta = 0.20$):
    \begin{equation}
    \tilde{\mathbf{w}}_1 = \mathbf{w}_0 - \eta \nabla f(\mathbf{w}_0) = \begin{bmatrix} 0.0 \\ 0.0 \end{bmatrix} - 0.20 \begin{bmatrix} -6.0 \\ -16.0 \end{bmatrix} = \begin{bmatrix} 0.0 + 1.20 \\ 0.0 + 3.20 \end{bmatrix} = \begin{bmatrix} 1.20 \\ 3.20 \end{bmatrix}
    \end{equation}
    Notice that $\tilde{\mathbf{w}}_1 = (1.20, 3.20)$ violates both upper bounds ($w_1 \le 1$ and $w_2 \le 1$).
    \item Orthogonal projection onto $\mathcal{C}$:
    \begin{equation}
    \mathbf{w}_1 = \Pi_{\mathcal{C}}(\tilde{\mathbf{w}}_1) = \begin{bmatrix} \min(1, \max(0, 1.20)) \\ \min(1, \max(0, 3.20)) \end{bmatrix} = \begin{bmatrix} 1.0 \\ 1.0 \end{bmatrix}
    \end{equation}
    \item Objective value after Iteration 1:
    \begin{equation}
    f(\mathbf{w}_1) = f(1.0, 1.0) = (1 - 3)^2 + 4(1 - 2)^2 = (-2)^2 + 4(-1)^2 = 4 + 4(1) = 8.0
    \end{equation}
    The loss plunges dramatically from $25.0 \to 8.0$ (a $68.0\%$ error reduction in a single iteration).
\end{itemize}

\textbf{Step 3: Iteration 2 Walkthrough}
\begin{itemize}
    \item Gradient computation at $\mathbf{w}_1 = [1.0, 1.0]^T$:
    \begin{equation}
    \nabla f(\mathbf{w}_1) = \begin{bmatrix} 2(1 - 3) \\ 8(1 - 2) \end{bmatrix} = \begin{bmatrix} 2(-2) \\ 8(-1) \end{bmatrix} = \begin{bmatrix} -4.0 \\ -8.0 \end{bmatrix}
    \end{equation}
    \item Intermediate unconstrained gradient step:
    \begin{equation}
    \tilde{\mathbf{w}}_2 = \mathbf{w}_1 - \eta \nabla f(\mathbf{w}_1) = \begin{bmatrix} 1.0 \\ 1.0 \end{bmatrix} - 0.20 \begin{bmatrix} -4.0 \\ -8.0 \end{bmatrix} = \begin{bmatrix} 1.0 + 0.80 \\ 1.0 + 1.60 \end{bmatrix} = \begin{bmatrix} 1.80 \\ 2.60 \end{bmatrix}
    \end{equation}
    \item Orthogonal projection onto $\mathcal{C}$:
    \begin{equation}
    \mathbf{w}_2 = \Pi_{\mathcal{C}}(\tilde{\mathbf{w}}_2) = \begin{bmatrix} \min(1, \max(0, 1.80)) \\ \min(1, \max(0, 2.60)) \end{bmatrix} = \begin{bmatrix} 1.0 \\ 1.0 \end{bmatrix}
    \end{equation}
    \item Convergence check:
    \begin{equation}
    \|\mathbf{w}_2 - \mathbf{w}_1\|_2 = \|[1.0, 1.0]^T - [1.0, 1.0]^T\|_2 = 0.0
    \end{equation}
    The iterate has locked onto the stationary fixed point $\mathbf{w}^* = (1.0, 1.0)$ in exactly one step!
\end{itemize}

\textbf{Step 4: Verification of First-Order Variational Optimality}
At the terminal iterate $\mathbf{w}^* = [1.0, 1.0]^T$, the objective gradient is $\nabla f(\mathbf{w}^*) = [-4.0, -8.0]^T$.
Testing the first-order variational inequality for \emph{any} arbitrary feasible point $\mathbf{z} \in [0, 1] \times [0, 1]$:
\begin{align}
\nabla f(\mathbf{w}^*)^T (\mathbf{z} - \mathbf{w}^*) &= \begin{bmatrix} -4.0 & -8.0 \end{bmatrix} \begin{bmatrix} z_1 - 1.0 \\ z_2 - 1.0 \end{bmatrix} \nonumber \\
&= -4.0(z_1 - 1.0) - 8.0(z_2 - 1.0) \nonumber \\
&= 4.0(1.0 - z_1) + 8.0(1.0 - z_2)
\end{align}
Because $\mathbf{z} \in \mathcal{C}$, we have $z_1 \le 1.0 \implies 1.0 - z_1 \ge 0$, and $z_2 \le 1.0 \implies 1.0 - z_2 \ge 0$.
Since both terms are strictly non-negative, their weighted sum is non-negative:
\begin{equation}
\nabla f(\mathbf{w}^*)^T (\mathbf{z} - \mathbf{w}^*) \ge 0 \quad \forall \mathbf{z} \in \mathcal{C}
\end{equation}
with equality if and only if $\mathbf{z} = \mathbf{w}^*$. This rigorously proves that $\mathbf{w}^* = (1.0, 1.0)$ is the unique global constrained minimizer of $f$ over the box $\mathcal{C}$.
\end{examplebox}

% ----------------------------------------------------------------------
% PART 6: EXPLICIT GOTCHAS & DIAGNOSTIC MISCONCEPTIONS
% ----------------------------------------------------------------------
\section{Explicit Gotchas \& Diagnostic Misconceptions}

\begin{trapbox}{Gotcha 1: Assuming the Union of Two Convex Sets is Convex}
A practitioner considers the union $S = B_1 \cup B_2$ of two disjoint convex norm balls $B_1 = \{\mathbf{x} \mid \|\mathbf{x} - \mathbf{c}_1\| \le 1\}$ and $B_2 = \{\mathbf{x} \mid \|\mathbf{x} - \mathbf{c}_2\| \le 1\}$ and assumes $S$ is convex.
\tcblower
\textbf{The Trap:} While the \textbf{intersection} of arbitrary convex sets is always convex, the \textbf{union} of convex sets is almost never convex. A straight line segment connecting a point in $B_1$ to a point in $B_2$ passes through the empty space between them, violating the convexity axiom. In optimization, unions of convex domains form non-convex feasible regions that create disconnected local minima.
\end{trapbox}

\begin{trapbox}{Gotcha 2: Believing Local Minima Can Differ from Global Minima in Convex Problems}
A student worries that gradient descent might get trapped in a suboptimal local minimum when minimizing a convex loss function over a convex set.
\tcblower
\textbf{The Trap:} For a convex functional defined on a convex set, \textbf{every local minimum is automatically a global minimum}. Furthermore, if the functional is strictly convex, the global minimizer is uniquely defined. The challenge of non-global local minima and spurious attractors exists exclusively on non-convex landscapes.
\end{trapbox}

\begin{trapbox}{Gotcha 3: Confusing Strict Convexity with Strong Convexity}
An analyst asserts that because $f(x) = x^4$ is strictly convex on $\mathbb{R}$, gradient descent will achieve linear convergence.
\tcblower
\textbf{The Trap:} $f(x) = x^4$ is strictly convex, but its second derivative $f''(x) = 12x^2$ vanishes at the minimizer $x = 0$. It is \textbf{not} $\mu$-strongly convex ($\mu = 0$). Strong convexity requires a uniform positive lower bound on curvature: $\nabla^2 f(\mathbf{x}) \succeq \mu I$ with $\mu > 0$. Without strong convexity, convergence near flat regions slows down to sublinear $\mathcal{O}(1/k)$.
\end{trapbox}

\begin{trapbox}{Gotcha 4: Forgetting Constraint Qualifications in Lagrange Systems}
A practitioner sets up a Lagrangian for the constraint $h(x, y) = x^2 + y^2 = 0$, but finds no Lagrange multipliers satisfying the stationarity equations.
\tcblower
\textbf{The Trap:} The Lagrange multiplier theorem requires a \textbf{constraint qualification}: the gradients of the active constraints must be linearly independent (LICQ). At the unique feasible point $(0, 0)$, $\nabla h(0, 0) = [2x, 2y]^T = [0, 0]^T$. The gradient vanishes, violating regularity. Always verify regularity ($\nabla h(\mathbf{w}^*) \ne \mathbf{0}$) before setting up the Lagrangian.
\end{trapbox}

\begin{trapbox}{Gotcha 5: The Boundary Curse: Projecting the Unconstrained Minimum}
When asked to minimize an objective $f(\mathbf{w})$ over a convex set $C$, an engineer computes the unconstrained minimum $\mathbf{w}_{\text{unc}}^* = \argmin_{\mathbf{w}} f(\mathbf{w})$ and projects it onto the boundary: $\hat{\mathbf{w}} = \Pi_C(\mathbf{w}_{\text{unc}}^*)$.
\tcblower
\textbf{The Trap:} The projection $\Pi_C(\mathbf{w}_{\text{unc}}^*)$ is \textbf{almost never the true constrained minimum}! Orthogonal projection minimizes Euclidean distance $\|\mathbf{w} - \mathbf{w}_{\text{unc}}^*\|_2^2$, treating all directions with equal isotropic weighting. If the objective function has anisotropic curvature (different eigenvalues in its Hessian), the true constrained optimum occurs at a point of \textbf{tangency} where $\nabla f(\mathbf{w}^*)$ is normal to the constraint boundary, not at the closest Euclidean point.
\end{trapbox}

\begin{trapbox}{Gotcha 6: The "Convex Loss on Non-Convex Parameterization" Illusion}
An ML engineer reasons: ``Mean Squared Error is quadratic and therefore strictly convex. Therefore, fitting a matrix factorization model $W \approx U V^T$ or a two-layer linear network $\hat{\mathbf{y}} = W_2 W_1 \mathbf{x}$ by minimizing squared reconstruction error has no spurious local minima or saddle points.''
\tcblower
\textbf{The Trap:} The composition of a convex function with a non-linear parameterization is \textbf{almost never convex}. 

While the loss functional $\phi(W) = \frac{1}{2}\|Y - W\|_F^2$ is strictly convex with respect to the matrix entries $W$, re-parameterizing $W = U V^T$ transforms the objective into:
\begin{equation}
f(U, V) = \frac{1}{2}\|Y - U V^T\|_F^2
\end{equation}
which is a degree-4 polynomial in the parameters $(U, V)$. Although $f(U, V)$ is \textbf{biconvex} (convex in $U$ for fixed $V$, and convex in $V$ for fixed $U$), it is jointly non-convex. It introduces saddle points at the origin $(U=O, V=O)$ where $\nabla f = \mathbf{0}$ but the Hessian is indefinite, as well as permutation symmetries and non-isolated manifolds of critical points. Never confuse the convexity of an outer loss with the convexity of the parameterized optimization problem.
\end{trapbox}

% ----------------------------------------------------------------------
% PART 7: COMPARATIVE SYNTHESIS & SUMMARY TABLE
% ----------------------------------------------------------------------
\section{Comparative Synthesis \& Summary Table}

Table~\ref{tab:convex_synthesis} summarizes the hierarchy of convexity classes for functionals.

\begin{table}[htbp]
\centering
\small
\begin{tabularx}{\textwidth}{lXXXX}
\toprule
\textbf{Convexity Class} & \textbf{Hessian Condition} & \textbf{First-Order Lower Bound} & \textbf{Minimizer Properties} & \textbf{Canonical Example} \\
\midrule
Convex & $\nabla^2 f(\mathbf{x}) \succeq 0$ & $f(\mathbf{y}) \ge f(\mathbf{x}) + \nabla f(\mathbf{x})^T(\mathbf{y} - \mathbf{x})$ & Any local minimum is a global minimum & $f(\mathbf{w}) = \|\mathbf{w}\|_1$ (Lasso loss) \\
\addlinespace
Strictly Convex & $\nabla^2 f(\mathbf{x}) \succ 0$ a.e. & Strict inequality for $\mathbf{y} \ne \mathbf{x}$ & Unique global minimizer (if it exists) & $f(x) = x^4$, $f(x) = e^x$ \\
\addlinespace
$\mu$-Strongly Convex & $\nabla^2 f(\mathbf{x}) \succeq \mu I$ ($\mu > 0$) & $f(\mathbf{y}) \ge f(\mathbf{x}) + \nabla f^T(\mathbf{y}-\mathbf{x}) + \frac{\mu}{2}\|\mathbf{y}-\mathbf{x}\|^2$ & Unique global minimizer; linear convergence & Ridge Regression: $\frac{1}{2}\|X\mathbf{w}-\mathbf{y}\|^2 + \frac{\lambda}{2}\|\mathbf{w}\|^2$ \\
\addlinespace
Non-Convex & Indefinite / Negative $\nabla^2 f$ & No global affine lower bound & Multiple local minima, saddle points & Neural networks, GMM negative log-likelihood \\
\bottomrule
\end{tabularx}
\caption{Comparative Synthesis: Hierarchy of Function Convexity Classes.}
\label{tab:convex_synthesis}
\end{table}

% ----------------------------------------------------------------------
% PART 8: PYTHON / SCIENTIFIC COMPUTING LAB
% ----------------------------------------------------------------------
\section{Python / Scientific Computing Lab}

The production script below implements Projected Gradient Descent (PGD) onto an $L_1$ ball constraint, demonstrating how projection enforces exact parameter sparsity and contrasting the result with naive unconstrained descent.

\begin{lstlisting}[language=Python, caption={Projected Gradient Descent onto an $L_1$ Ball in Python.}]
import numpy as np

def project_onto_l1_ball(v: np.ndarray, s: float) -> np.ndarray:
    """Computes Euclidean projection of vector v onto L1 ball of radius s: ||w||_1 <= s."""
    if np.sum(np.abs(v)) <= s:
        return v
    # Sort absolute values descending
    u = np.sort(np.abs(v))[::-1]
    cssv = np.cumsum(u)
    # Find active threshold coordinate
    rho = np.nonzero(u * np.arange(1, len(v) + 1) > (cssv - s))[0][-1]
    theta = (cssv[rho] - s) / (rho + 1.0)
    w = np.sign(v) * np.maximum(np.abs(v) - theta, 0.0)
    return w

def run_projected_gradient_descent(X: np.ndarray, y: np.ndarray, s_radius: float, lr: float, n_iters: int):
    """Minimizes 0.5 * ||Xw - y||^2 subject to ||w||_1 <= s_radius via PGD."""
    N, d = X.shape
    w = np.zeros(d)
    losses = []

    for k in range(n_iters):
        grad = X.T @ (X @ w - y)
        w_unconstrained = w - lr * grad
        w = project_onto_l1_ball(w_unconstrained, s_radius)
        loss = 0.5 * float(np.sum((X @ w - y) ** 2))
        losses.append(loss)

    return w, losses

def main():
    rng = np.random.RandomState(42)
    N, d = 100, 10
    X = rng.randn(N, d)
    
    # Ground truth sparse weights (only 3 active components)
    w_true = np.zeros(d)
    w_true[1] = 2.5
    w_true[4] = -1.8
    w_true[7] = 3.0
    y = X @ w_true + 0.1 * rng.randn(N)

    # Maximum Lipschitz smoothness constant L = lambda_max(X^T X)
    L = float(np.linalg.eigvalsh(X.T @ X).max())
    lr = 1.0 / L

    s_radius = 5.0
    w_pgd, losses = run_projected_gradient_descent(X, y, s_radius, lr, n_iters=200)
    w_ols = np.linalg.lstsq(X, y, rcond=None)[0]

    print("=================================================================")
    print(" PROJECTED GRADIENT DESCENT: L1 BALL SPARSITY LAB")
    print("=================================================================")
    print(f"L1 Ball Radius Constraint:      ||w||_1 <= {s_radius:.2f}")
    print(f"Unconstrained OLS Weights Norm:  ||w_ols||_1 = {np.sum(np.abs(w_ols)):.4f}")
    print(f"Projected PGD Weights Norm:      ||w_pgd||_1 = {np.sum(np.abs(w_pgd)):.4f}")
    print(f"Exact Zeros in OLS Solution:     {np.sum(np.isclose(w_ols, 0))} / {d}")
    print(f"Exact Zeros in PGD Solution:     {np.sum(np.isclose(w_pgd, 0))} / {d}")
    print(f"Final Constrained Loss:          {losses[-1]:.4f}")
    print("-" * 65)

    print(f"{'Index':<8}{'True Weight':<16}{'OLS (Unconstrained)':<22}{'PGD (L1 Constrained)':<20}")
    print("-" * 65)
    for i in range(d):
        print(f"{i:<8}{w_true[i]:<16.2f}{w_ols[i]:<22.4f}{w_pgd[i]:<20.4f}")

if __name__ == "__main__":
    main()
\end{lstlisting}

% ----------------------------------------------------------------------
% PART 9: COMPREHENSIVE PRACTICE PROBLEMS
% ----------------------------------------------------------------------
\section{Comprehensive Practice Problems}

\subsection*{Part A: Analytical \& Mathematical Problems}

\begin{enumerate}[label=\textbf{P\arabic*.}]
    \item \textbf{Analytical Lagrange Multipliers with Linear Constraint:}
    Minimize the quadratic objective $f(x, y) = 2x^2 + y^2$ subject to the linear equality constraint $x + y = 3$.
    \begin{enumerate}[label=(\alph*)]
        \item Formulate the Lagrangian $\mathcal{L}(x, y, \nu)$.
        \item Solve the stationarity system $\nabla_{x, y, \nu} \mathcal{L} = \mathbf{0}$.
        \item Compute the optimal coordinates $(x^*, y^*)$ and optimal multiplier $\nu^*$.
        \item Verify that $\nabla f(x^*, y^*)$ is collinear with $\nabla h(x^*, y^*)$.
    \end{enumerate}

    \item \textbf{Maximum Entropy Probability Distribution:}
    Let $\mathbf{p} = [p_1, \dots, p_K]^T$ be a discrete probability vector. Maximize the Shannon entropy $H(\mathbf{p}) = -\sum_{i=1}^K p_i \log p_i$ subject to the simplex constraint $\sum_{i=1}^K p_i = 1$.
    Using Lagrange multipliers, prove that the uniform distribution $p_i^* = \frac{1}{K}$ for all $i$ is the unique global maximizer.

    \item \textbf{Convexity Verification (Second-Order Conditions):}
    Consider the log-sum-exp function $f(\mathbf{x}) = \log\big( \sum_{i=1}^d e^{x_i} \big)$.
    \begin{enumerate}[label=(\alph*)]
        \item Compute the gradient $\nabla f(\mathbf{x})$ and show it is a probability vector (softmax).
        \item Compute the Hessian matrix $\nabla^2 f(\mathbf{x})$.
        \item Prove that $\nabla^2 f(\mathbf{x}) \succeq 0$ for all $\mathbf{x} \in \mathbb{R}^d$, establishing that log-sum-exp is convex.
    \end{enumerate}

    \item \textbf{Projection onto a Halfspace:}
    Let $\mathcal{H}_- = \{\mathbf{x} \in \mathbb{R}^d \mid \mathbf{a}^T \mathbf{x} \le b\}$ with $\|\mathbf{a}\|_2 = 1$.
    Prove that the Euclidean projection of a point $\mathbf{x}_0 \in \mathbb{R}^d$ onto $\mathcal{H}_-$ is given by:
    \begin{equation}
    \Pi_{\mathcal{H}_-}(\mathbf{x}_0) = \mathbf{x}_0 - \max(0, \mathbf{a}^T \mathbf{x}_0 - b) \mathbf{a}
    \end{equation}

    \item \textbf{Jensen's Inequality Application:}
    Using the convexity of $f(t) = -\log t$ on $(0, \infty)$, prove the Arithmetic Mean - Geometric Mean (AM-GM) inequality: for any positive real numbers $a_1, a_2, \dots, a_n > 0$,
    \begin{equation}
    \left( \prod_{i=1}^n a_i \right)^{1/n} \le \frac{1}{n} \sum_{i=1}^n a_i
    \end{equation}

    \item \textbf{Variational Characterization of Projection \& Firm Non-Expansiveness:}
    Let $C \subset \mathbb{R}^d$ be a non-empty closed convex set.
    \begin{enumerate}[label=(\alph*)]
        \item Let $\mathbf{w} \in \mathbb{R}^d$. Prove that $\mathbf{z}^* = \Pi_C(\mathbf{w})$ if and only if:
        \begin{equation*}
        \langle \mathbf{w} - \mathbf{z}^*, \mathbf{x} - \mathbf{z}^* \rangle \le 0 \quad \forall \mathbf{x} \in C
        \end{equation*}
        \item Using part (a), prove that the projection operator is \emph{firmly non-expansive}:
        \begin{equation*}
        \|\Pi_C(\mathbf{u}) - \Pi_C(\mathbf{v})\|_2^2 + \|(I - \Pi_C)(\mathbf{u}) - (I - \Pi_C)(\mathbf{v})\|_2^2 \le \|\mathbf{u} - \mathbf{v}\|_2^2 \quad \forall \mathbf{u}, \mathbf{v} \in \mathbb{R}^d
        \end{equation*}
        \item Deduce that $\|\Pi_C(\mathbf{u}) - \Pi_C(\mathbf{v})\|_2 \le \|\mathbf{u} - \mathbf{v}\|_2$, establishing that $\Pi_C$ is a 1-Lipschitz contraction mapping.
    \end{enumerate}

    \item \textbf{Lagrange Multiplier Sensitivity \& Shadow Price Derivation:}
    Consider the perturbed constrained optimization problem $f^*(\mathbf{c}) = \min_{\mathbf{w}} f(\mathbf{w})$ subject to $h_j(\mathbf{w}) = c_j$ for $j = 1, \dots, p$, where $f$ and all $h_j$ are smooth.
    \begin{enumerate}[label=(\alph*)]
        \item Set up the Lagrangian $\mathcal{L}(\mathbf{w}, \boldsymbol{\nu}; \mathbf{c}) = f(\mathbf{w}) + \sum_{j=1}^p \nu_j (h_j(\mathbf{w}) - c_j)$.
        \item Assuming the optimal trajectory $\mathbf{w}^*(\mathbf{c})$ is differentiable and satisfies the LICQ condition at $\mathbf{c} = \mathbf{0}$, prove using the multivariable chain rule that:
        \begin{equation*}
        \frac{\partial f^*(\mathbf{c})}{\partial c_k} = -\nu_k^*(\mathbf{c})
        \end{equation*}
        \item An automated trading algorithm allocates a total capital budget $B = \$100{,}000$ across assets by minimizing portfolio variance subject to $\sum w_i = B$. If the optimal Lagrange multiplier is $\nu^* = -0.0042$, calculate the exact marginal change in minimal portfolio variance if the capital budget is increased by $\$1{,}000$.
    \end{enumerate}
\end{enumerate}

\subsection*{Part B: Diagnostic Multiple-Choice Questions}

\begin{enumerate}[label=\textbf{MCQ\arabic*.}]
    \item Which of the following sets is NOT convex?
    \begin{enumerate}[label=(\alph*)]
        \item The Euclidean unit ball $\{\mathbf{x} \in \mathbb{R}^d \mid \|\mathbf{x}\|_2 \le 1\}$.
        \item The positive semidefinite cone $S^d_+ = \{A \in \mathbb{R}^{d \times d} \mid A = A^T, A \succeq 0\}$.
        \item The hollow spherical shell $\{\mathbf{x} \in \mathbb{R}^d \mid \|\mathbf{x}\|_2 = 1\}$.
        \item The halfspace $\{\mathbf{x} \in \mathbb{R}^d \mid \mathbf{a}^T \mathbf{x} \le b\}$.
    \end{enumerate}

    \item If $f: \mathbb{R}^d \to \mathbb{R}$ is convex and continuously differentiable, what is the geometric relationship between the graph of $f$ and the tangent hyperplane at any point $\mathbf{x}_0$?
    \begin{enumerate}[label=(\alph*)]
        \item The tangent hyperplane lies strictly above the graph of $f$.
        \item The tangent hyperplane is a global under-estimator that lies everywhere below or touches the graph of $f$.
        \item The tangent hyperplane intersects the graph at infinitely many isolated points.
        \item The tangent hyperplane has a slope equal to the Hessian determinant.
    \end{enumerate}
\end{enumerate}
